% ================= PORTADA =================
\thispagestyle{empty}
\begin{center}
{\small UNIVERSIDAD NACIONAL AUTÓNOMA DE MÉXICO · Posgrado en Ciencia e Ingeniería de la Computación}\\[1.2cm]
{\color{azul}\rule{\linewidth}{1pt}}\\[0.4cm]
{\Huge\bfseries\color{azul} Guía de estudio}\\[0.3cm]
{\LARGE Programación Avanzada — Primer examen parcial}\\[0.3cm]
{\large Prof. Gustavo Márquez Flores · Semestre 2026-1}\\[0.4cm]
{\color{azul}\rule{\linewidth}{1pt}}\\[0.8cm]
\ifsol
{\Large\bfseries\color{verde} CUADERNO DE SOLUCIONES}\\[0.2cm]
{\normalsize Cada problema aparece con su enunciado y su solución desarrollada paso a paso.}
\else
{\Large\bfseries\color{azul} CUADERNO DE PROBLEMAS}\\[0.2cm]
{\normalsize Resuelve primero sin mirar; después compara con el cuaderno de soluciones (misma numeración).}
\fi
\end{center}
\vspace{0.8cm}

\begin{formulario}[Contenido]
\begin{enumerate}[label=\textbf{\arabic*.},leftmargin=*]
\item Fundamentos: lenguajes, gramáticas, paradigmas, programación estructurada y modular, métricas de módulos.
\item Divide y vencerás: esquema general, búsqueda binaria, QuickSort, multiplicación de enteros grandes.
\item Backtracking: una, todas y la mejor solución; ocho reinas; cuadros mágicos y construcción por patrones.
\item Lógica de Hoare básica: aserciones, fortaleza, corrección parcial/total, asignación, consecuencia, conjunción.
\item Composición secuencial e invariantes de bloque.
\item Condicionales (\textbf{if} con y sin \textbf{else}).
\item Ciclos: inducción, invariante de ciclo, salida y terminación.
\item Recursión doble: función 91 de McCarthy, inducción modificada y función de Ackermann.
\item Programación lógica: cláusulas de Horn, resolución, SLD, unificación, backtracking y corte en Prolog.
\item Dos exámenes simulados (formato de 5--6 preguntas, como el del año pasado).
\end{enumerate}
Los problemas marcados con $\bigstar$ son \emph{literalmente} las preguntas del examen del año pasado.
\end{formulario}

\newpage
% ================= CÓMO ES EL EXAMEN =================
\section*{Cómo es el examen del profesor}
\addcontentsline{toc}{section}{Cómo es el examen}

El primer parcial del año pasado tuvo \textbf{cinco preguntas} que combinan demostración formal, teoría breve y pseudocódigo:

\begin{center}
\begin{tabular}{@{}c p{10.4cm} p{3.9cm}@{}}
\toprule
\# & Pregunta & Tipo \\ \midrule
1 & Demostrar que $\hoare{a<b}{\kw{if}\ (a>b)\ \kw{then}\ m\asg a\ \kw{else}\ m\asg b}{m=b}$ es válida. & Demostración (Hoare) \\
2 & ¿Qué métricas de módulos existen para medir su eficiencia? & Teoría breve \\
3 & Escribir en seudocódigo el algoritmo de \emph{Divide y vencerás}. & Pseudocódigo \\
4 & Demostrar que el invariante de \texttt{int j=9; for(int i=0;i<10;i++) j--;} es $i+j=9$. & Demostración (inducción) \\
5 & Describir el algoritmo de \emph{Backtracking} para buscar la \emph{mejor} solución. & Pseudocódigo + teoría \\
\bottomrule
\end{tabular}
\end{center}

\textbf{Qué se deduce:} el profesor pide (i) demostraciones con reglas de Hoare escritas completas, (ii) invariantes de ciclos probados por inducción, (iii) los esquemas de sus hojas \emph{``Algoritmo BackTracking y Divide-Venceras''} en su misma notación y (iv) definiciones cortas. Como este año el curso ya avanzó a la función 91, Ackermann y Prolog, es razonable esperar al menos una pregunta de esos temas.

\begin{consejo}[Lección del examen corregido]
En la pregunta 1 el alumno supuso $a>b$ y trató de ``forzar'' $m=b$; el profesor anotó \emph{``revisar planteamiento''}. La clave era \textbf{reconocer que la rama \textbf{then} es inalcanzable}: $a<b \land a>b\equiv\bot$, y una terna con precondición $\bot$ es válida por vacuidad. Siempre escribe explícitamente $P$, $B$, $Q$, las dos ramas y la regla que usas.
\end{consejo}

\subsection*{Plantillas que conviene memorizar}
\begin{formulario}[Plantilla: demostrar una terna con condicional]
\begin{enumerate}[leftmargin=*]
\item Identificar $P$ (precondición), $B$ (guarda) y $Q$ (postcondición).
\item \textbf{Rama verdadera:} con precondición $P\land B$, calcular $wp(C_1,Q)$ por sustitución y probar $P\land B\Rightarrow wp(C_1,Q)$ (regla de consecuencia). Si $P\land B\equiv\bot$, la rama es inalcanzable y la terna vale por vacuidad.
\item \textbf{Rama falsa:} igual con $P\land\neg B$ y $C_2$ (si no hay \textbf{else}: probar $P\land\neg B\Rightarrow Q$).
\item Concluir con la regla condicional. Sin ciclos ni recursión, el código termina: corrección total.
\end{enumerate}
\end{formulario}

\begin{formulario}[Plantilla: demostrar un invariante de ciclo]
\begin{enumerate}[leftmargin=*]
\item Traducir el \texttt{for} a \texttt{while} y nombrar valores tras $n$ iteraciones: $x_n$.
\item \textbf{Caso base} ($n=0$): el invariante $I$ vale antes de la primera iteración.
\item \textbf{Hipótesis de inducción:} $I$ vale tras $n$ iteraciones.
\item \textbf{Paso inductivo:} si la guarda $B$ es cierta y se ejecuta el cuerpo, $I$ vale tras $n+1$ iteraciones (equivalente: $\hoare{I\land B}{\text{cuerpo}}{I}$).
\item \textbf{Salida:} $I\land\neg B\Rightarrow Q$. \quad \textbf{Terminación:} variante entera $\ge0$ que decrece.
\end{enumerate}
\end{formulario}

\newpage
% ================= FORMULARIO =================
\section*{Formulario de repaso}
\addcontentsline{toc}{section}{Formulario}

\begin{formulario}[Lógica de Hoare]
\renewcommand{\arraystretch}{1.35}
\begin{tabular}{@{}p{4.3cm} p{11.3cm}@{}}
Terna & $\hoare{P}{C}{Q}$: si $C$ inicia en un estado que satisface $P$ \emph{y termina}, el estado final satisface $Q$. \\[2pt]
Corrección & total $=$ parcial $+$ terminación. \\[2pt]
Fortaleza & $R$ es más fuerte que $S$ si $R\Rightarrow S$.\quad $\bot$ (ningún estado) es la más fuerte; $\T\equiv\{\ \}$ (todos) la más débil. \\[2pt]
Asignación & $\hoare{\sust{Q}{E}{V}}{V\asg E}{Q}$ \quad (sustituir $V$ por $E$ en $Q$; añadir condiciones de dominio, p.\,ej. divisor $\neq0$). \\[2pt]
Consecuencia & $\dfrac{P_1\Rightarrow P\quad \hoare{P}{C}{Q}}{\hoare{P_1}{C}{Q}}$ \qquad $\dfrac{\hoare{P}{C}{Q}\quad Q\Rightarrow Q_1}{\hoare{P}{C}{Q_1}}$ \\[10pt]
Conjunción / disyunción & $\dfrac{\hoare{P_1}{C}{Q_1}\quad\hoare{P_2}{C}{Q_2}}{\hoare{P_1\land P_2}{C}{Q_1\land Q_2}}$ \qquad $\dfrac{\hoare{P_1}{C}{Q_1}\quad\hoare{P_2}{C}{Q_2}}{\hoare{P_1\lor P_2}{C}{Q_1\lor Q_2}}$ \\[10pt]
Composición & $\dfrac{\hoare{P}{C_1}{R}\quad\hoare{R}{C_2}{Q}}{\hoare{P}{C_1;C_2}{Q}}$ \quad ($R$: aserción intermedia; se calcula hacia atrás). \\[10pt]
\textbf{if} sin \textbf{else} & $\dfrac{\hoare{P\land B}{C_1}{Q}\quad (P\land\neg B)\Rightarrow Q}{\hoare{P}{\kw{if}\ B\ \kw{then}\ C_1}{Q}}$ \\[10pt]
\textbf{if} con \textbf{else} & $\dfrac{\hoare{P\land B}{C_1}{Q}\quad \hoare{P\land\neg B}{C_2}{Q}}{\hoare{P}{\kw{if}\ B\ \kw{then}\ C_1\ \kw{else}\ C_2}{Q}}$ \\[10pt]
\textbf{while} & $P\Rightarrow I$;\ \ $\hoare{I\land B}{C}{I}$;\ \ $I\land\neg B\Rightarrow Q$;\ \ variante $V\in\mathbb N$ que decrece estrictamente. \\[2pt]
Invariante de bloque & $\hoare{I}{C}{I}$: se conserva la \emph{relación}, no los valores.
\end{tabular}
\end{formulario}

\begin{formulario}[Inducción, McCarthy y Ackermann]
\textbf{Inducción ordinaria:} $P(k)$ y $\forall n\ge k\,(P(n)\Rightarrow P(n+1))$ $\Rightarrow$ $\forall n\ge k\ P(n)$.\\
\textbf{Inducción modificada (hacia abajo):} $P(k)$ y, para todo $n<k$, [$P(m)$ para todo $m$ con $n<m\le k$] $\Rightarrow P(n)$ $\;\Longrightarrow\;$ $\forall n\le k\ P(n)$.\\[3pt]
\textbf{Función 91:} $f(x)=x-10$ si $x>100$; $f(x)=f(f(x+11))$ si $x\le100$. \ Teorema: $f(x)=91$ para $x\le 100$.\\[3pt]
\textbf{Ackermann:} $A(0,y)=y+1$;\ $A(x+1,0)=A(x,1)$;\ $A(x+1,y+1)=A(x,A(x+1,y))$.\\
$A(1,z)=z+2$,\quad $A(2,z)=2z+3$,\quad $A(3,z)=2^{z+3}-3$,\quad $A(4,z)=\underbrace{2^{2^{\cdot^{\cdot^{2}}}}}_{z+3\text{ doses}}-3$ \ ($A(4,0)=13$, $A(4,1)=65533$).
\end{formulario}

\begin{formulario}[Técnicas de diseño]
\textbf{Divide y vencerás:} caso base simple; descomponer en $k$ subproblemas \emph{semejantes y menores}; resolver recursivamente; \emph{combinar}.\\
Recurrencias típicas: $T(n)=T(n/2)+\Theta(1)\Rightarrow\Theta(\log n)$ (búsqueda binaria); $2T(n/2)+\Theta(1)\Rightarrow\Theta(n)$ (máximo); $2T(n/2)+\Theta(n)\Rightarrow\Theta(n\log n)$ (mergesort, QuickSort balanceado); $T(n-1)+\Theta(n)\Rightarrow\Theta(n^2)$ (QuickSort peor caso); $4T(n/2)+\Theta(n)\Rightarrow\Theta(n^2)$; $3T(n/2)+\Theta(n)\Rightarrow\Theta(n^{\log_2 3})$ (Karatsuba).\\[3pt]
\textbf{Backtracking:} árbol de decisiones recorrido en profundidad; \emph{elegir – validar (poda) – anotar – avanzar – deshacer}. \emph{Una}: se detiene con el primer éxito. \emph{Todas}: registra y sigue. \emph{Mejor}: compara con la mejor guardada y sigue (con cota $\to$ \emph{branch and bound}).\\[3pt]
\textbf{Ocho reinas:} una reina por fila; columna $j$, diagonales $j-k$ y $j+k$. 92 soluciones (12 fundamentales). $n=4$: 2 soluciones.\\
\textbf{Cuadro mágico:} $M_n=\dfrac{n(n^2+1)}{2}$ ($M_3=15$, $M_4=34$, $M_5=65$). Impar: método siamés (arriba-derecha; si está ocupada, abajo). $n\equiv0\pmod 4$: llenar $1..n^2$ y cambiar $x\mapsto n^2+1-x$ fuera de las diagonales de cada bloque $4\times4$.
\end{formulario}

\begin{formulario}[Prolog y resolución]
Regla \ \texttt{a :- b1, ..., bm.} \ $\equiv$ \ $b_1\land\dots\land b_m\to a$ \ $\equiv$ \ $a\lor\neg b_1\lor\dots\lor\neg b_m$ (cláusula de Horn definida).\\
Meta \ \texttt{?- a1, ..., an.} \ $\equiv$ \ $\leftarrow A_1,\dots,A_n$ \ $\equiv$ \ $\neg A_1\lor\dots\lor\neg A_n$.\\
Resolución: $\dfrac{C\lor L\qquad D\lor\neg L'}{(C\lor D)\theta}$,\ \ $\theta=\mathrm{UMG}(L,L')$. Refutación: $S\cup\{\neg F\}\vdash\square\;\Rightarrow\;S\models F$.\\
Paso SLD: $\leftarrow A_1,\dots,A_k,\dots,A_n$ con regla $A\leftarrow B_1..B_m$ y $\theta=\mathrm{UMG}(A,A_k)$ da $\leftarrow(A_1..A_{k-1},B_1..B_m,A_{k+1}..A_n)\theta$. Éxito: meta vacía $\square$.\\
Sintaxis: \texttt{,} conjunción; \texttt{;} disyunción (u otra respuesta en consola); \texttt{\textbackslash+} negación como falla; \texttt{!} corte; \texttt{is} evalúa aritmética; \texttt{=} solo unifica; \texttt{=:=} compara valores; Mayúscula inicial o \texttt{\_} = variable; minúscula = átomo; sin espacio entre predicado y paréntesis.
\end{formulario}

\newpage
% =====================================================================
\section{Fundamentos y técnicas de programación}
% =====================================================================

\begin{prob}
¿Qué significa que un programa sea \emph{correcto}? Explica por qué ``correcto'' y ``eficiente'' no son sinónimos y menciona tres criterios con los que suele juzgarse un programa.
\end{prob}
\begin{sol}
Un programa es \textbf{correcto respecto de una especificación} (precondición $P$ y postcondición $Q$) si, para toda entrada que satisface $P$, produce una salida que satisface $Q$ (corrección parcial) y además siempre termina (corrección total). La corrección es \emph{relativa} a la especificación: no existe ``correcto en abstracto''.

La eficiencia mide \emph{cuántos recursos} usa el programa (tiempo, memoria). Un programa puede ser correcto y lentísimo (p.\,ej.\ ordenar probando todas las permutaciones) o rápido e incorrecto.

Criterios (preguntas del ``tour'' del curso): ¿está bien escrito y es legible?, ¿es óptimo en \textbf{tiempo} y \textbf{memoria}?, ¿usa eficientemente el lenguaje?, ¿el algoritmo es eficiente?, ¿se puede \textbf{demostrar} correcto?
\end{sol}

\begin{prob}
Define \emph{algoritmo} y \emph{programa}. ¿Qué expresa la fórmula de Niklaus Wirth?
\end{prob}
\begin{sol}
\textbf{Algoritmo:} método de solución con un número finito de pasos bien definidos; puede recibir datos, produce al menos un resultado y termina en tiempo finito. \textbf{Programa:} un algoritmo expresado en un lenguaje de programación (código fuente sintácticamente válido) para que una computadora lo ejecute.

Wirth: \emph{Algoritmos $+$ Estructuras de datos $=$ Programas}. Resolver un problema no es solo describir pasos: también hay que decidir cómo representar y organizar la información sobre la que operan.
\end{sol}

\begin{prob}
Sea $G=(V,T,S,P)$ con $T=\{a,b\}$, $N=\{S\}$ y producciones $S\to aSb\mid ab$.
\begin{enumerate}[label=(\alph*)]
\item Explica qué es cada componente de $G$.
\item Deriva la cadena $aaabbb$. ¿Cuál es $L(G)$?
\item ¿De qué tipo es $G$ en la jerarquía de Chomsky? ¿Y la gramática $S\to aS\mid b$?
\end{enumerate}
\end{prob}
\begin{sol}
(a) $V$: vocabulario finito; $T\subseteq V$: terminales; $N=V\setminus T$: no terminales; $S\in N$: símbolo inicial; $P$: reglas de producción.

(b) $S\Rightarrow aSb\Rightarrow aaSbb\Rightarrow aaabbb$. \quad $L(G)=\{a^nb^n\mid n\ge1\}$.

(c) Cada producción tiene un único no terminal a la izquierda: \textbf{tipo 2, libre de contexto}. No es regular (tipo 3): $a^nb^n$ exige ``contar'', cosa que un autómata finito no puede. La gramática $S\to aS\mid b$ tiene producciones $A\to aB$ o $A\to a$: \textbf{tipo 3, regular}; genera $L=\{a^nb\mid n\ge0\}$.

Tabla: tipo 0 irrestricta, tipo 1 sensible al contexto, tipo 2 libre de contexto, tipo 3 regular. La sintaxis de casi todos los lenguajes de programación es libre de contexto; comprobaciones como ``variable declarada'' o ``tipos compatibles'' requieren análisis adicional.
\end{sol}

\begin{prob}
Clasifica por paradigma (imperativo / declarativo funcional / declarativo lógico) y por forma de ejecución típica: C, Pascal, Java, Haskell, Lisp, Prolog, Python. ¿Por qué la clasificación compilado/interpretado no es absoluta?
\end{prob}
\begin{sol}
\begin{tabular}{@{}lll@{}}
\toprule Lenguaje & Paradigma & Ejecución típica\\\midrule
C, Pascal & imperativo (estructurado) & compilado\\
Java & imperativo, orientado a objetos & compilado a bytecode + JVM (interpreta/JIT)\\
Python & imperativo / OO (multiparadigma) & interpretado\\
Haskell & declarativo funcional & compilado (GHC) o interpretado (GHCi)\\
Lisp & declarativo funcional & interpretado o compilado\\
Prolog & declarativo lógico & interpretado (SWI) o compilado (GNU Prolog)\\\bottomrule
\end{tabular}

Imperativo: describe \emph{cómo} cambiar el estado. Declarativo: describe \emph{qué} relación o resultado se quiere. La forma de ejecución depende de la \emph{implementación}, no del lenguaje: un mismo lenguaje puede tener compilador, intérprete y JIT.
\end{sol}

\begin{prob}
Enuncia el teorema de Böhm--Jacopini y da un ejemplo en pseudocódigo de cada estructura.
\end{prob}
\begin{sol}
Todo algoritmo (programa con un punto de entrada y uno de salida) puede expresarse combinando solo tres estructuras de control:
\begin{enumerate}
\item \textbf{Secuencia:} \texttt{x := 1; y := x + 2}.
\item \textbf{Selección:} \texttt{if x > 0 then y := x else y := -x}.
\item \textbf{Iteración:} \texttt{while i < n do i := i + 1}.
\end{enumerate}
Es la base de la programación estructurada (sin \texttt{goto}). El teorema no garantiza legibilidad: los nombres, la modularidad y la organización siguen importando.
\end{sol}

\begin{prob}[\examen]
¿Qué métricas de módulos existen para medir su eficiencia (y calidad)?
\end{prob}
\begin{sol}
\textbf{Métricas de diseño modular vistas en clase:}
\begin{itemize}
\item \textbf{Cohesión:} medida de qué tan relacionadas están las responsabilidades \emph{internas} de un módulo. Se busca \textbf{alta} cohesión (el módulo hace una sola cosa bien definida), lo que facilita nombrarlo, entenderlo, probarlo y reutilizarlo.
\item \textbf{Acoplamiento:} medida de la \emph{dependencia} de un módulo respecto de otros. Se busca \textbf{bajo} acoplamiento, para poder cambiar un módulo sin afectar a los demás.
\end{itemize}
Objetivo de diseño: \textbf{alta cohesión y bajo acoplamiento}.

\textbf{Métricas de eficiencia del módulo:}
\begin{itemize}
\item \textbf{Tiempo de ejecución} / complejidad temporal (p.\,ej.\ $O(n\log n)$).
\item \textbf{Uso de memoria} / complejidad espacial.
\end{itemize}
\textbf{Métrica de corrección:} si el módulo es correcto (corrección parcial) y si termina (corrección total).

\emph{Complemento (útil si piden ejemplos, no desarrollado en clase):} tipos de cohesión de Yourdon--Constantine (funcional, secuencial, comunicacional, procedimental, temporal, lógica, coincidental; de mejor a peor) y de acoplamiento (de datos, por estampado, de control, externo, común, de contenido; de mejor a peor); complejidad ciclomática de McCabe $V(G)=E-N+2$; líneas de código; \emph{fan-in}/\emph{fan-out}.
\end{sol}

\begin{prob}
Las diapositivas advierten que ``alta cohesión $\Rightarrow$ alto acoplamiento'' no es la meta. Explica la idea y describe cómo MVC reparte responsabilidades entre módulos.
\end{prob}
\begin{sol}
Dividir un sistema en módulos muy cohesivos obliga a que se comuniquen, lo que \emph{podría} aumentar el acoplamiento. La meta es lograr alta cohesión \emph{sin} dependencias innecesarias: módulos que hablan mediante interfaces pequeñas y bien definidas.

\textbf{MVC:} \emph{Controlador} recibe la solicitud y coordina (p.\,ej.\ servlet); \emph{Modelo} contiene datos y lógica del dominio (JavaBeans/EJB, acceso a datos); \emph{Vista} presenta el resultado (JSP). Cada parte tiene una responsabilidad (alta cohesión) y la vista puede cambiar sin tocar la lógica del dominio (bajo acoplamiento).
\end{sol}

% =====================================================================
\section{Divide y vencerás}
% =====================================================================

\begin{prob}[\examen]
Escribe en seudocódigo el algoritmo de \emph{Divide y vencerás} para resolver un problema y explica las condiciones que debe cumplir el problema para aplicarlo.
\end{prob}
\begin{sol}
Esquema del profesor:
\begin{lstlisting}
Funcion Divide_y_Venceras( P : Problema ): Solución
{ Devuelve la solución del problema P }
Inicio
  Si P es suficientemente pequeño o simple
    Entonces devolver( solucion_simple( P ) )
  Si_No
    descomponer P en k partes( P1, P2, ..., Pk );
    Para( i = 1; i <= k; i++ )
      Si = Divide_y_Venceras( Pi );
    devolver( combinar( S1, S2, ..., Sk ) );
Fin
\end{lstlisting}
\textbf{Condiciones:}
\begin{enumerate}
\item $P$ se puede descomponer en subproblemas $P_1,\dots,P_k$ \textbf{semejantes} al original y \textbf{de menor tamaño}.
\item Existe un \textbf{caso base} que se resuelve directamente (\texttt{solucion\_simple}).
\item Las subsoluciones se pueden \textbf{combinar} eficientemente.
\item Cada llamada reduce una medida (p.\,ej.\ el tamaño), lo que garantiza la \textbf{terminación}.
\end{enumerate}
Costo: $T(n)=\sum_i T(n_i)+D(n)+C(n)$ (dividir $+$ combinar); con $a$ subproblemas de tamaño $n/b$: $T(n)=aT(n/b)+f(n)$.

Ejemplos: QuickSort, búsqueda binaria, mergesort, multiplicación de enteros grandes.

\emph{Ojo:} en la hoja del ``tour'' aparece \texttt{combinar(P1,...,Pk)}; lo correcto es combinar las \emph{soluciones} $S_1,\dots,S_k$.
\end{sol}

\begin{prob}
Escribe la búsqueda binaria recursiva sobre un arreglo ordenado $A[izq..der]$. Traza la búsqueda de $X=23$ y de $X=10$ en $A=[3,8,12,17,23,31,40,52]$ (índices $1..8$). ¿Por qué las llamadas usan $m-1$ y $m+1$? ¿Cuál es su complejidad?
\end{prob}
\begin{sol}
\begin{lstlisting}
Funcion BusBin( A, X, izq, der ): entero
Inicio
  Si izq > der entonces devolver( 0 )        // caso vacío: no está
  m = (izq + der) div 2
  Si X = A[m] entonces devolver( m )          // caso base
  Si X < A[m] entonces devolver( BusBin(A, X, izq, m-1) )
  Si_No devolver( BusBin(A, X, m+1, der) )
Fin
\end{lstlisting}
\begin{tabular}{@{}l|l@{}}
$X=23$ & $X=10$\\\hline
$[1,8]$, $m=4$, $A[4]=17<23\to[5,8]$ & $[1,8]$, $m=4$, $A[4]=17>10\to[1,3]$\\
$[5,8]$, $m=6$, $A[6]=31>23\to[5,5]$ & $[1,3]$, $m=2$, $A[2]=8<10\to[3,3]$\\
$[5,5]$, $m=5$, $A[5]=23$ \ $\Rightarrow$ devuelve 5 & $[3,3]$, $m=3$, $A[3]=12>10\to[3,2]$: vacío $\Rightarrow$ 0
\end{tabular}

Usar $m\pm1$ excluye el elemento ya comparado; así el intervalo \emph{se reduce estrictamente} en cada llamada y la recursión termina (con $m$ en lugar de $m-1$ podría repetirse el mismo intervalo indefinidamente). Es divide y vencerás con un solo subproblema y combinación trivial: $T(n)=T(n/2)+\Theta(1)=\Theta(\log n)$. Requiere arreglo \textbf{ordenado}.
\end{sol}

\begin{prob}
Diseña con divide y vencerás una función \texttt{Maximo(A, l, r)} que devuelva el mayor elemento de $A[l..r]$. Plantea la recurrencia y resuélvela.
\end{prob}
\begin{sol}
\begin{lstlisting}
Funcion Maximo( A, l, r ): valor
Inicio
  Si l = r entonces devolver( A[l] )            // caso simple
  m = (l + r) div 2                             // descomponer en 2
  M1 = Maximo( A, l, m )
  M2 = Maximo( A, m+1, r )
  Si M1 >= M2 entonces devolver( M1 )           // combinar
  Si_No devolver( M2 )
Fin
\end{lstlisting}
$T(1)=\Theta(1)$, $T(n)=2T(n/2)+\Theta(1)$. Al desarrollar: $T(n)=2^kT(n/2^k)+(2^k-1)c$; con $2^k=n$ queda $T(n)=\Theta(n)$ (exactamente $n-1$ comparaciones). Terminación: $l\le m<r$, así que ambos subintervalos son estrictamente menores.
\end{sol}

\begin{prob}
Considera el procedimiento \texttt{Sort(l, r)} de \texttt{Qsort.PAS}:
\begin{lstlisting}
i := l; j := r; x := a[(l+r) div 2];
repeat
  while a[i] < x do i := i + 1;
  while x < a[j] do j := j - 1;
  if i <= j then begin
    y := a[i]; a[i] := a[j]; a[j] := y;  i := i + 1; j := j - 1;
  end;
until i > j;
if l < j then Sort(l, j);
if i < r then Sort(i, r);
\end{lstlisting}
Traza la primera partición para $a=[5,3,8,1,9,2,7]$ (índices $1..7$) y di qué llamadas recursivas se generan. ¿Qué enseña este ejemplo sobre el pivote? Da la complejidad en el mejor y peor caso.
\end{prob}
\begin{sol}
Pivote: $x=a[(1+7)\ \mathrm{div}\ 2]=a[4]=1$.
\begin{itemize}
\item $i$: $a[1]=5\not<1\Rightarrow i=1$. \quad $j$: $1<7,\,1<2,\,1<9$ retrocede; $1\not<a[4]=1\Rightarrow j=4$.
\item $i\le j$: intercambiar $a[1]\leftrightarrow a[4]$: $[1,3,8,5,9,2,7]$; $i=2$, $j=3$.
\item $i$: $a[2]=3\not<1\Rightarrow i=2$. \quad $j$: $1<8$, $1<3$ retrocede; $a[1]=1\Rightarrow j=1$.
\item $i=2>j=1$: no hay intercambio y termina el \texttt{repeat}.
\end{itemize}
Llamadas: $l<j$? $1<1$ no. $\;i<r$? $2<7$ sí $\Rightarrow$ \texttt{Sort(2,7)}. Continúa: \texttt{Sort(2,7)} con pivote 5 $\to$ $[1,3,2,5,9,8,7]$, luego \texttt{Sort(2,3)} y \texttt{Sort(5,7)}; resultado $[1,2,3,5,7,8,9]$.

\textbf{Lección:} tomar el elemento \emph{central por posición} no garantiza una partición equilibrada: aquí el pivote fue el mínimo y quedaron regiones de tamaño 0 y 6. Complejidad: $T(n)=T(k)+T(n-k-1)+\Theta(n)$. Balanceado: $\Theta(n\log n)$ (caso típico/esperado). Muy desbalanceado: $T(n)=T(n-1)+\Theta(n)=\Theta(n^2)$. Memoria de pila: $O(\log n)$ típico, $O(n)$ peor caso.

\emph{Invariante de la partición:} lo recorrido a la izquierda de $i$ es $\le x$ y lo recorrido a la derecha de $j$ es $\ge x$.
\end{sol}

\begin{prob}
Multiplica $u=1234$ y $v=5678$ con el esquema de divide y vencerás de clase ($s=2$, $B=10^2$). Da la recurrencia del esquema y explica la mejora de Karatsuba.
\end{prob}
\begin{sol}
$u=Bw+x$ con $w=12$, $x=34$; \ $v=By+z$ con $y=56$, $z=78$. Por distributividad
\[uv=B^2\,wy+B\,(wz+xy)+xz.\]
$wy=672$, $wz=936$, $xy=1904$, $xz=2652$.
\[uv=672\cdot10^4+(936+1904)\cdot10^2+2652=6\,720\,000+284\,000+2\,652=\mathbf{7\,006\,652}.\]
Recurrencia: 4 multiplicaciones de la mitad de dígitos $+$ sumas/desplazamientos lineales: $T(n)=4T(n/2)+\Theta(n)=\Theta(n^2)$ — no mejora al método escolar.

\textbf{Karatsuba:} calcula $wy$, $xz$ y $p=(w+x)(y+z)$; entonces $wz+xy=p-wy-xz$. Solo 3 multiplicaciones: $T(n)=3T(n/2)+\Theta(n)=\Theta(n^{\log_2 3})\approx\Theta(n^{1.585})$. Aquí: $p=46\cdot134=6164$ y $6164-672-2652=2840=936+1904$. \checkmark
\end{sol}

\begin{prob}
Escribe \emph{mergesort} siguiendo el esquema del profesor, señalando qué parte es ``descomponer'' y cuál ``combinar''. Compáralo con QuickSort: ¿dónde hace cada uno el trabajo pesado?
\end{prob}
\begin{sol}
\begin{lstlisting}
Funcion MergeSort( A, l, r )
Inicio
  Si l >= r entonces devolver                     // 0 o 1 elemento: ya ordenado
  m = (l + r) div 2                               // descomponer (trivial)
  MergeSort( A, l, m );  MergeSort( A, m+1, r )   // resolver
  Mezclar( A, l, m, r )                           // combinar: Theta(n)
Fin
\end{lstlisting}
$T(n)=2T(n/2)+\Theta(n)=\Theta(n\log n)$ \emph{siempre}, con $\Theta(n)$ de memoria auxiliar. En \textbf{mergesort} dividir es trivial y el trabajo está en \textbf{combinar} (mezclar). En \textbf{QuickSort} el trabajo está en \textbf{dividir} (partición alrededor del pivote) y combinar es trivial (las regiones ya quedan en su lugar, \emph{in-place}).
\end{sol}

\begin{prob}
Responde brevemente: (a) ¿Todo algoritmo recursivo es divide y vencerás? (b) ¿Qué diferencia hay entre caso base y caso vacío? (c) ¿Cuál es la diferencia esencial entre divide y vencerás y backtracking?
\end{prob}
\begin{sol}
(a) No. La recursión es un \emph{mecanismo de implementación}; divide y vencerás es una \emph{estrategia de diseño} (dividir en subproblemas semejantes, resolver, combinar). Ej.: la función 91 es recursiva pero no divide ni combina.

(b) El caso base da una solución directa (p.\,ej.\ $A[m]=X$); el caso vacío indica que ya no hay candidatos ($izq>der$).

(c) Divide y vencerás \emph{necesita todas} las subsoluciones y las combina. Backtracking explora un árbol de decisiones, \emph{abandona} ramas inviables (poda) y deshace decisiones para probar alternativas; busca una, todas o la mejor solución.
\end{sol}

% =====================================================================
\section{Backtracking, ocho reinas y cuadros mágicos}
% =====================================================================

\begin{prob}[\examen]
Describe el algoritmo de \emph{Backtracking} cuando se trata de buscar la \textbf{mejor} solución de un problema. Indica en qué difiere de las versiones que buscan \emph{una} y \emph{todas} las soluciones.
\end{prob}
\begin{sol}
Backtracking construye la solución por \textbf{etapas} $i$; en cada una prueba las opciones disponibles, \textbf{anota} solo las \emph{aceptables} (poda), avanza a la etapa $i+1$ y, al regresar, \textbf{cancela la anotación} para probar la siguiente opción (vuelta atrás). Recorre en profundidad el árbol de soluciones parciales.

Versión \emph{mejor solución} (notación de la hoja del profesor):
\begin{lstlisting}
Funcion Backtracking( Etapa i, solucionMejor ) devuelve: boolean
Inicio
  Éxito = falso;
  IniciarOpciones( i, GrupoOpciones o );
  Repetir
    SeleccionarNuevaOpcion( o, Opcion n );
    Si ( Aceptable( n ) ) entonces
      AnotarOpcion( i, n );
      Si SolucionCompleta( i ) entonces
        Éxito = verdadero;
        Si EsMejorSolucion( solucionActual, solucionMejor ) entonces
          ApuntarSolucionActual();      // solucionMejor := solucionActual
        finsi;
      Sino
        Éxito = Backtracking( i+1, solucionMejor );
        Si Éxito = falso entonces cancelamosAnotacion( i, n ); finsi;
      Finsi;
    Finsi;
  Hasta ( NoQuedanOpciones( o ) );       // NO se detiene al primer éxito
  Retorna Éxito;
Fin
\end{lstlisting}
\textbf{Diferencias:}
\begin{itemize}
\item \textbf{Una solución:} el ciclo termina con \texttt{Hasta (Éxito = verdadero) o NoQuedanOpciones(o)}: se detiene en la primera hoja válida.
\item \textbf{Todas:} el ciclo es \texttt{Hasta NoQuedanOpciones(o)}; cada solución completa se guarda con \texttt{ApuntarSoluciónCompleta(i, solucionesCompletas)} y se sigue explorando.
\item \textbf{Mejor:} también explora todo el árbol, pero en vez de guardar todas compara cada solución completa con la mejor encontrada (\texttt{EsMejorSolucion}) y conserva solo la mejor. Requiere un \textbf{criterio de calidad}. Si además se podan ramas cuya cota no puede superar a la mejor actual, se obtiene \emph{ramificación y acotamiento} (branch and bound).
\end{itemize}
\emph{Detalle fino:} para seguir explorando después de registrar una solución hay que deshacer la anotación \emph{siempre} (no solo cuando falla la llamada); en muchas implementaciones se escribe \texttt{cancelamosAnotacion(i,n)} al final de cada intento. Costo sin poda efectiva: $O(b^d)$ nodos (ramificación $b$, profundidad $d$).
\end{sol}

\begin{prob}
Problema de las 4 reinas (una reina por fila, filas $k=0..3$, columnas $j=1..4$).
\begin{enumerate}[label=(\alph*)]
\item ¿Por qué colocar una reina por fila elimina la restricción de filas sin comprobarla?
\item ¿Qué identifican $j-k$ y $j+k$? ¿Se atacan $(k,j)=(0,2)$ y $(2,4)$? ¿Y $(1,4)$ y $(3,2)$?
\item Traza el árbol de búsqueda en profundidad (probando columnas de 1 a 4) hasta la primera solución.
\item ¿Cuántas soluciones tiene el problema de 4 reinas? ¿Y el de 8? ¿Qué son las 12 soluciones fundamentales?
\end{enumerate}
\end{prob}
\begin{sol}
(a) Cada nivel $k$ de la recursión decide solo la columna de la fila $k$; es imposible poner dos reinas en la misma fila, así que esa restricción se cumple \emph{por construcción}. La solución se representa como $\mathit{sol}[k]=j$.

(b) $j-k$ es constante en cada diagonal ``$\searrow$'' ($45^\circ$ en el código) y $j+k$ en cada diagonal ``$\swarrow$'' ($135^\circ$). $(0,2)$ y $(2,4)$: $2-0=2=4-2$ $\Rightarrow$ \textbf{se atacan}. $(1,4)$ y $(3,2)$: $4+1=5=2+3$ $\Rightarrow$ \textbf{se atacan}.

(c) Notación: $F_k=j$.
\begin{itemize}
\item $F_0=1$. $F_1$: 1 (columna), 2 (diag.\ $j-k$: $1=1$) \texttimes; $F_1=3$ \checkmark. $F_2$: 1 col.\ \texttimes, 2 ($j+k=4$ vs.\ $(1,3)$: $4$) \texttimes, 3 col.\ \texttimes, 4 ($j-k=2$ vs.\ $(1,3)$: $2$) \texttimes\ $\Rightarrow$ \emph{retroceder}.
\item $F_1=4$ \checkmark. $F_2=2$ \checkmark. $F_3$: 1,2,4 columnas ocupadas, 3 ($j-k=0$ vs.\ $(2,2)$) \texttimes\ $\Rightarrow$ retroceder. $F_2=3$ ($j-k=1$ vs.\ $(0,1)$) \texttimes; 4 col.\ \texttimes\ $\Rightarrow$ retroceder hasta $F_0$.
\item $F_0=2$. $F_1$: 1 ($j+k=2$) \texttimes, 2 col.\ \texttimes, 3 ($j-k=2$) \texttimes, $F_1=4$ \checkmark. $F_2=1$ \checkmark. $F_3$: 1 col., 2 col., $F_3=3$ \checkmark\ ($j-k=0$, $j+k=6$ libres).
\end{itemize}
Primera solución: $\mathit{sol}=[2,4,1,3]$.

(d) 4 reinas: \textbf{2} soluciones, $[2,4,1,3]$ y $[3,1,4,2]$ (una es reflejo de la otra). 8 reinas: \textbf{92}; agrupando las que se obtienen por rotaciones y reflexiones del tablero quedan \textbf{12 fundamentales}. El programa de clase imprime 92 porque no elimina simetrías.
\end{sol}

\begin{prob}
En el programa \texttt{Ocho reinas.c} visto en clase:
\begin{lstlisting}
for (int j = 1; j <= NREINAS; j++)
  if (!contiene(col, j) && !contiene(diag45, j - k) && !contiene(diag135, j + k)) {
    sol[k] = j;
    col.push_back(j); diag45.push_back(j - k); diag135.push_back(j + k);
    reinas(k + 1, col, diag45, diag135);
    col.pop_back(); diag45.pop_back(); diag135.pop_back();
  }
\end{lstlisting}
(a) Relaciona cada línea con las operaciones del esquema de backtracking del profesor. (b) ¿Qué representa \texttt{k == NREINAS}? (c) ¿Por qué no hace falta ``borrar'' \texttt{sol[k]}?
\end{prob}
\begin{sol}
(a) \texttt{for j} $=$ \texttt{IniciarOpciones}/\texttt{SeleccionarNuevaOpcion}; el \texttt{if} con \texttt{contiene} $=$ \texttt{Aceptable(n)} (poda); \texttt{sol[k]=j} y los \texttt{push\_back} $=$ \texttt{AnotarOpcion}; \texttt{reinas(k+1,...)} $=$ \texttt{Backtracking(i+1)}; los \texttt{pop\_back} $=$ \texttt{cancelamosAnotacion} (deshacer). El ciclo recorre todas las columnas: versión \emph{todas las soluciones}.

(b) Que ya se colocaron las 8 reinas (filas $0..7$): \texttt{SolucionCompleta}; se imprime y se regresa.

(c) \texttt{sol[k]} se sobrescribe antes de volver a imprimir una solución completa. (Los vectores se pasan por valor, así que los \texttt{pop\_back} expresan el retroceso aunque no son indispensables para el nivel llamador.)
\end{sol}

\begin{prob}[\simil]
Mochila $0/1$: objetos con pesos $w=(2,3,4,5)$ y valores $v=(3,4,5,7)$, capacidad $7$. Se quiere el subconjunto de mayor valor.
\begin{enumerate}[label=(\alph*)]
\item Formúlalo como backtracking de \emph{mejor solución}: ¿qué es una etapa, una opción, \texttt{Aceptable}, \texttt{SolucionCompleta} y \texttt{EsMejorSolucion}?
\item Encuentra la mejor solución y muestra al menos una rama podada.
\end{enumerate}
\end{prob}
\begin{sol}
(a) Etapa $i$: decidir sobre el objeto $i$ ($i=1..4$). Opciones: \{incluir, no incluir\}. \texttt{Aceptable}: el peso acumulado no excede 7 (poda por factibilidad). \texttt{SolucionCompleta}: $i=4$ (se decidió sobre todos). \texttt{EsMejorSolucion}: valor actual $>$ mejor valor guardado. Poda adicional opcional (cota): si valor actual $+$ suma de valores restantes $\le$ mejor, no vale la pena seguir.

(b) Soluciones factibles relevantes: $\{1,2\}$: peso 5, valor 7; $\{1,3\}$: 6, 8; $\{1,4\}$: 7, \textbf{10}; $\{2,3\}$: 7, 9; $\{4\}$: 5, 7. Mejor: $\{1,4\}$ (pesos $2+5$) con valor \textbf{10}.\\
Ramas podadas por factibilidad: tras incluir 1 y 2 (peso 5), incluir el 3 daría peso 9 $>7$ $\Rightarrow$ no \texttt{Aceptable}; lo mismo con el 4 (peso 10). Tras $\{1,3\}$ (peso 6), incluir el 4 da 11: podado.\\
Poda por cota (branch and bound): cuando ya se tiene la mejor $=10$, la rama ``no incluir 1, no incluir 2, no incluir 3'' puede lograr a lo más $v_4=7<10$, así que se descarta sin explorarla.
\end{sol}

\begin{prob}
Cuadros mágicos.
\begin{enumerate}[label=(\alph*)]
\item Deduce la constante mágica $M_n$ y calcula $M_3, M_4, M_5, M_6$.
\item Construye el cuadro $3\times3$ y el $5\times5$ con el método siamés.
\item Construye el $4\times4$ con el método de complementos.
\item ¿Por qué el método siamés no sirve para $n=4$ ni el de complementos para $n=6$?
\end{enumerate}
\end{prob}
\begin{sol}
(a) La suma de todas las casillas es $1+\dots+n^2=\frac{n^2(n^2+1)}{2}$; hay $n$ filas con la misma suma $M_n$: $nM_n=\frac{n^2(n^2+1)}{2}\Rightarrow M_n=\frac{n(n^2+1)}{2}$. $M_3=15$, $M_4=34$, $M_5=65$, $M_6=111$.

(b) Siamés: 1 en el centro de la fila superior; el siguiente va arriba-derecha (con ``envoltura'' $(f-1)\bmod n$, $(c+1)\bmod n$); si esa casilla está ocupada, se coloca \emph{debajo} de la actual.
\[
\begin{array}{|c|c|c|}\hline 8&1&6\\\hline 3&5&7\\\hline 4&9&2\\\hline\end{array}
\qquad
\begin{array}{|c|c|c|c|c|}\hline 17&24&1&8&15\\\hline 23&5&7&14&16\\\hline 4&6&13&20&22\\\hline 10&12&19&21&3\\\hline 11&18&25&2&9\\\hline\end{array}
\]
Todas las filas, columnas y diagonales suman 15 y 65, respectivamente.

(c) Llenar $1..16$ por filas; conservar las casillas de las diagonales ($i\bmod4=j\bmod4$ o $i\bmod 4+j\bmod4=3$) y sustituir las demás por $17-x$:
\[
\begin{array}{|c|c|c|c|}\hline 1&2&3&4\\\hline 5&6&7&8\\\hline 9&10&11&12\\\hline 13&14&15&16\\\hline\end{array}
\;\longrightarrow\;
\begin{array}{|c|c|c|c|}\hline 1&15&14&4\\\hline 12&6&7&9\\\hline 8&10&11&5\\\hline 13&3&2&16\\\hline\end{array}
\quad(\text{todas suman }34)
\]

(d) El siamés depende de que $n$ sea impar (hay columna central y el recorrido diagonal cubre todas las casillas sin conflicto); el de complementos necesita bloques $4\times4$ completos ($n\equiv0\pmod4$). $n=6$ es ``simplemente par'' y requiere otra construcción (p.\,ej.\ LUX/Strachey).
\end{sol}

\begin{prob}
Formula la búsqueda de \emph{todos} los cuadros mágicos $3\times3$ como backtracking: estado, opciones, condición de aceptación (poda) y solución completa. ¿Por qué para $n\ge5$ es preferible la construcción por patrones si solo se quiere uno?
\end{prob}
\begin{sol}
\textbf{Estado:} matriz parcialmente llena y conjunto de números usados. \textbf{Etapa:} siguiente casilla vacía (por filas). \textbf{Opciones:} números de $1..9$ no usados. \textbf{Aceptable (poda):} ninguna suma parcial de fila/columna/diagonal supera 15, y toda línea que se complete suma exactamente 15. \textbf{Solución completa:} las 9 casillas llenas (y todas las líneas suman 15). Versión \emph{todas}: registrar y seguir. Para $3\times3$ hay 8 soluciones (una fundamental con sus rotaciones y reflexiones).

El espacio crece como $(n^2)!$: 880 cuadros de orden 4 y 275\,305\,224 de orden 5. Un patrón (siamés, complementos) produce \emph{uno} directamente en $\Theta(n^2)$, sin explorar árbol alguno. Enumerar y construir son problemas distintos.
\end{sol}

% =====================================================================
\section{Lógica de Hoare: aserciones, asignación y consecuencia}
% =====================================================================

\begin{prob}
Define: aserción, estado, $s\models P$, terna de Hoare, corrección parcial y corrección total. ¿Por qué en un fragmento sin ciclos ni recursión basta la corrección parcial?
\end{prob}
\begin{sol}
\textbf{Aserción:} fórmula lógica sobre las variables del programa (p.\,ej.\ $x=3\land y\le z$). \textbf{Estado:} asignación de valores a las variables. $s\models P$: el estado $s$ satisface $P$. \textbf{Terna} $\hoare{P}{C}{Q}$: si $C$ empieza en un estado que satisface $P$ y termina, el estado final satisface $Q$. \textbf{Parcial:} la terna es válida (``si termina, el resultado es correcto''). \textbf{Total:} parcial $+$ se demuestra que $C$ termina para toda entrada que cumple $P$.

Un fragmento sin ciclos ni recursión ejecuta un número finito de instrucciones primitivas que terminan; la terminación es inmediata y parcial $\Rightarrow$ total.
\end{sol}

\begin{prob}
Ordena de la más fuerte a la más débil: $\;x\ge0,\ \ \T,\ \ x>10,\ \ \F,\ \ x=15,\ \ x>0$. ¿Qué significa la aserción vacía \{\ \}? ¿Es cierto que $y\neq0\Rightarrow y=4$?
\end{prob}
\begin{sol}
$\F\;\Rightarrow\;x=15\;\Rightarrow\;x>10\;\Rightarrow\;x>0\;\Rightarrow\;x\ge0\;\Rightarrow\;\T$.
Más fuerte $=$ admite menos estados. $\F$ no admite ninguno (la más fuerte); \{\ \} $\equiv\T$ admite todos (la más débil): $\T\land P\equiv P$, $\T\lor P\equiv\T$. La implicación no es simétrica: $y=4\Rightarrow y\ne0$, pero $y\ne0\not\Rightarrow y=4$ (contraejemplo $y=7$).
\end{sol}

\begin{prob}
Calcula la precondición más débil $P$ en cada caso:
\[
\begin{array}{ll}
\text{(a) } \hoare{P}{x\asg x+5}{x>12} & \text{(b) } \hoare{P}{y\asg 3x-1}{y\le 8}\\[2pt]
\text{(c) } \hoare{P}{x\asg x\cdot x}{x<9} & \text{(d) } \hoare{P}{z\asg z/2}{z=5}\\[2pt]
\text{(e) } \hoare{P}{a\asg b}{a=b} & \text{(f) } \hoare{P}{x\asg 10/(y-2)}{x>0}
\end{array}
\]
\end{prob}
\begin{sol}
Axioma: $P=\sust{Q}{E}{V}$ (sustituir la variable asignada por la expresión), más condiciones de dominio.
\begin{enumerate}[label=(\alph*)]
\item $x+5>12\iff \mathbf{x>7}$.
\item $3x-1\le8\iff\mathbf{x\le3}$.
\item $x^2<9\iff\mathbf{-3<x<3}$.
\item $z/2=5\iff\mathbf{z=10}$.
\item $b=b\iff\mathbf{\T}$: no hace falta ninguna condición; $\hoare{\ }{a\asg b}{a=b}$.
\item $\frac{10}{y-2}>0$ y dominio $y\ne2$: el cociente es positivo si y solo si $y-2>0$: $\mathbf{y>2}$.
\end{enumerate}
\end{sol}

\begin{prob}
¿Son válidas? Justifica con el axioma de asignación y la regla de consecuencia, o da un contraejemplo.
\[
\begin{array}{ll}
\text{(a) }\hoare{x=5}{x\asg x+1}{x>5} & \text{(b) }\hoare{x>0}{x\asg x-1}{x>0}\\[2pt]
\text{(c) }\hoare{\T}{y\asg x\cdot x}{y\ge0} & \text{(d) }\hoare{x\ge0}{x\asg x-1}{x\ge-1}
\end{array}
\]
\end{prob}
\begin{sol}
\begin{enumerate}[label=(\alph*)]
\item $wp=x+1>5\iff x>4$. Como $x=5\Rightarrow x>4$, por \textbf{fortalecimiento de la precondición} es \textbf{válida}.
\item $wp=x-1>0\iff x>1$. Pero $x>0\not\Rightarrow x>1$. Contraejemplo: $x=1$ termina con $x=0$. \textbf{No válida}.
\item $wp=x^2\ge0\iff\T$. \textbf{Válida} para cualquier estado.
\item $wp=x-1\ge-1\iff x\ge0$: la precondición dada es exactamente la $wp$. \textbf{Válida}.
\end{enumerate}
\end{sol}

\begin{prob}
(a) A partir de $\hoare{y\neq0}{x\asg 1/y}{x=1/y}$ deduce $\hoare{y=4}{x\asg1/y}{x=1/y}$ escribiendo la regla usada.
(b) A partir de $\hoare{\T}{max\asg b}{max=b}$ deduce $\hoare{\T}{max\asg b}{max\ge b}$.
(c) Explica con conjuntos de estados por qué se puede \emph{fortalecer} la precondición y \emph{debilitar} la postcondición, pero no al revés.
\end{prob}
\begin{sol}
(a) Fortalecimiento de la precondición:
$\dfrac{y=4\Rightarrow y\ne0\qquad\hoare{y\ne0}{x\asg1/y}{x=1/y}}{\hoare{y=4}{x\asg1/y}{x=1/y}}$.

(b) Debilitamiento de la postcondición:
$\dfrac{\hoare{\T}{max\asg b}{max=b}\qquad max=b\Rightarrow max\ge b}{\hoare{\T}{max\asg b}{max\ge b}}$.

(c) Si el código funciona para \emph{todos} los estados de $P$, funciona para cualquier subconjunto $P_1\subseteq P$ (más fuerte). Si el resultado cae en el conjunto $Q$, también cae en cualquier superconjunto $Q_1\supseteq Q$ (más débil). Al revés no: ampliar la precondición mete estados no verificados (en (a), $y=0$), y restringir la postcondición puede excluir resultados reales.
\end{sol}

\begin{prob}
Usa la regla de conjunción (y consecuencia) para demostrar $\hoare{i=i_0\land i_0\ge3}{i\asg i+2}{i\ge5}$, partiendo de las ternas $\hoare{i=i_0}{i\asg i+2}{i=i_0+2}$ y $\hoare{i_0\ge3}{i\asg i+2}{i_0\ge3}$. ¿Por qué la segunda terna es válida?
\end{prob}
\begin{sol}
La segunda es válida por el axioma de asignación: $i_0$ es una variable lógica (el valor inicial) que no aparece asignada, así que $\sust{(i_0\ge3)}{i+2}{i}=i_0\ge3$.
Conjunción:
$\dfrac{\hoare{i=i_0}{i\asg i+2}{i=i_0+2}\qquad\hoare{i_0\ge3}{i\asg i+2}{i_0\ge3}}{\hoare{i=i_0\land i_0\ge3}{i\asg i+2}{i=i_0+2\land i_0\ge3}}$.
Como $i=i_0+2\land i_0\ge3\Rightarrow i\ge5$, por debilitamiento de la postcondición: $\hoare{i=i_0\land i_0\ge3}{i\asg i+2}{i\ge5}$. \checkmark

(Directo con el axioma: $wp=i+2\ge5\iff i\ge3$.)
\end{sol}

% =====================================================================
\section{Composición secuencial e invariantes de bloque}
% =====================================================================

\begin{prob}
Demuestra $\hoare{\ }{c\asg a+b;\ c\asg c/2}{c=(a+b)/2}$ (ejercicio del ``tour'' del curso).
\end{prob}
\begin{sol}
Hacia atrás desde $Q: c=\frac{a+b}{2}$.
\begin{enumerate}
\item Por $c\asg c/2$: $\sust{Q}{c/2}{c}$: $\frac c2=\frac{a+b}2\iff c=a+b$. \quad (1) $\hoare{c=a+b}{c\asg c/2}{c=\tfrac{a+b}2}$.
\item Por $c\asg a+b$: $\sust{(c=a+b)}{a+b}{c}$: $a+b=a+b\iff\T$. \quad (2) $\hoare{\T}{c\asg a+b}{c=a+b}$.
\item Composición de (2) y (1) con aserción intermedia $R: c=a+b$:
$\dfrac{\hoare{\T}{c\asg a+b}{c=a+b}\quad\hoare{c=a+b}{c\asg c/2}{c=\frac{a+b}2}}{\hoare{\T}{c\asg a+b;\ c\asg c/2}{c=\frac{a+b}2}}$.
\end{enumerate}
La precondición \{\ \} $\equiv\T$: funciona para todos los valores de $a$ y $b$.
\end{sol}

\begin{prob}
Encuentra la precondición más débil: (a) $\hoare{P}{x\asg x+1;\ y\asg 2x}{y>10}$; \ (b) $\hoare{P}{y\asg x-3;\ z\asg y\cdot y}{z>16}$.
\end{prob}
\begin{sol}
(a) Por $y\asg2x$: $2x>10\iff x>5$. Por $x\asg x+1$: $x+1>5\iff x>4$. \ $P: x>4$.

(b) Por $z\asg y\cdot y$: $y^2>16\iff y>4\lor y<-4$. Por $y\asg x-3$: $x-3>4\lor x-3<-4\iff x>7\lor x<-1$. \ $P: x>7\lor x<-1$.
\end{sol}

\begin{prob}
Demuestra que el intercambio \emph{sin variable auxiliar} es correcto:
$\hoare{a=A\land b=B}{a\asg a+b;\ b\asg a-b;\ a\asg a-b}{a=B\land b=A}$.
\end{prob}
\begin{sol}
$A,B$ son constantes lógicas (valores iniciales). Hacia atrás:
\begin{enumerate}
\item Por $a\asg a-b$: $(a-b=B)\land(b=A)$. \hfill (3)
\item Por $b\asg a-b$: sustituir $b$ por $a-b$ en (3): $(a-(a-b)=B)\land(a-b=A)\iff b=B\land a-b=A$. \hfill (2)
\item Por $a\asg a+b$: sustituir $a$ por $a+b$ en (2): $b=B\land (a+b)-b=A\iff b=B\land a=A$. \hfill (1)
\end{enumerate}
(1) coincide con la precondición. Componiendo las tres ternas:
$\hoare{a=A\land b=B}{a\asg a+b}{b=B\land a-b=A}$, $\hoare{b=B\land a-b=A}{b\asg a-b}{a-b=B\land b=A}$, $\hoare{a-b=B\land b=A}{a\asg a-b}{a=B\land b=A}$. \checkmark\ (Con enteros de máquina podría haber desbordamiento en $a+b$.)
\end{sol}

\begin{prob}
Demuestra que $I: s=i^2$ es invariante del bloque $\;i\asg i+1;\ s\asg s+2i-1$. Comprueba con $i=3,s=9$.
\end{prob}
\begin{sol}
Hay que probar $\hoare{s=i^2}{i\asg i+1;\ s\asg s+2i-1}{s=i^2}$.
\begin{enumerate}
\item Por $s\asg s+2i-1$: $s+2i-1=i^2\iff s=(i-1)^2$. \quad $\hoare{s=(i-1)^2}{s\asg s+2i-1}{s=i^2}$.
\item Por $i\asg i+1$: $s=((i+1)-1)^2=i^2$. \quad $\hoare{s=i^2}{i\asg i+1}{s=(i-1)^2}$.
\item Composición $\Rightarrow\hoare{s=i^2}{i\asg i+1;\ s\asg s+2i-1}{s=i^2}$. \checkmark
\end{enumerate}
Con valores: $i_1=4$, $s_1=9+2\cdot4-1=16=4^2$. El invariante no dice que los valores no cambien, sino que se conserva la \emph{relación}.
\end{sol}

\begin{prob}
Demuestra $\hoare{\ }{s\asg1;\ s\asg s\cdot r+2;\ s\asg s\cdot r+3}{s=r^2+2r+3}$ (evaluación de Horner). Observa dónde es necesaria la regla de consecuencia.
\end{prob}
\begin{sol}
\begin{enumerate}
\item Por $s\asg s\cdot r+3$: $wp=(s\cdot r+3=r^2+2r+3)\iff s\cdot r=r(r+2)$. Esta $wp$ es incómoda (si $r=0$ cualquier $s$ sirve). Usamos una aserción más fuerte: $R_2: s=r+2$, que \emph{implica} la $wp$. Por consecuencia: $\hoare{s=r+2}{s\asg s\cdot r+3}{s=r^2+2r+3}$.
\item Por $s\asg s\cdot r+2$ con postcondición $R_2$: $wp=(s\cdot r+2=r+2)\iff s\cdot r=r$. Tomamos $R_1: s=1$, que la implica: $\hoare{s=1}{s\asg s\cdot r+2}{s=r+2}$.
\item Por $s\asg1$: $1=1\iff\T$: $\hoare{\T}{s\asg1}{s=1}$.
\item Composición de las tres. \checkmark
\end{enumerate}
\textbf{Lección:} la aserción intermedia no tiene que ser la $wp$ exacta; basta una condición más fuerte que la implique (regla de consecuencia).
\end{sol}

% =====================================================================
\section{Condicionales}
% =====================================================================

\begin{prob}[\examen]
Demostrar que la siguiente terna de Hoare es válida:
\[\hoare{a<b}{\kw{if}\ (a>b)\ \kw{then}\ m\asg a\ \kw{else}\ m\asg b}{m=b}.\]
\end{prob}
\begin{sol}
Identificamos $P: a<b$, \ $B: a>b$, \ $Q: m=b$. Por la regla del \textbf{if} con \textbf{else} basta probar
(1) $\hoare{P\land B}{m\asg a}{Q}$ y (2) $\hoare{P\land\neg B}{m\asg b}{Q}$.

\textbf{(1) Rama verdadera.} $P\land B\equiv a<b\land a>b\equiv\F$ (no existe estado con $a<b$ y $a>b$ a la vez). Formalmente: por asignación, $\hoare{a=b}{m\asg a}{m=b}$ (sustituir $m$ por $a$ en $m=b$). Como $\F\Rightarrow a=b$ (de lo falso se sigue cualquier cosa), por fortalecimiento de la precondición:
\[\hoare{a<b\land a>b}{m\asg a}{m=b}\quad\text{es válida (por vacuidad): la rama es \emph{inalcanzable}.}\]

\textbf{(2) Rama falsa.} $P\land\neg B\equiv a<b\land a\le b\equiv a<b$. Por asignación: $\sust{(m=b)}{b}{m}\equiv b=b\equiv\T$, es decir, $\hoare{\T}{m\asg b}{m=b}$. Como $a<b\Rightarrow\T$, por consecuencia:
\[\hoare{a<b\land\neg(a>b)}{m\asg b}{m=b}.\]

\textbf{Conclusión.} Por la regla condicional,
\[\dfrac{\hoare{a<b\land a>b}{m\asg a}{m=b}\qquad\hoare{a<b\land\neg(a>b)}{m\asg b}{m=b}}{\hoare{a<b}{\kw{if}\ (a>b)\ \kw{then}\ m\asg a\ \kw{else}\ m\asg b}{m=b}}\]
Como no hay ciclos ni recursión, el programa termina: la terna es \textbf{totalmente correcta}. \ $\blacksquare$

\emph{Intuición:} si $a<b$, la guarda $a>b$ es falsa, siempre se ejecuta $m\asg b$ y por tanto $m=b$.
\end{sol}

\begin{prob}
Demuestra $\hoare{\ }{\kw{if}\ a>b\ \kw{then}\ m\asg a\ \kw{else}\ m\asg b}{m\ge a\land m\ge b}$.
\end{prob}
\begin{sol}
$P=\T$, $B: a>b$, $Q: m\ge a\land m\ge b$.

\textbf{Rama verdadera} ($a>b$): $\sust{Q}{a}{m}=a\ge a\land a\ge b\equiv a\ge b$. Como $a>b\Rightarrow a\ge b$: $\hoare{a>b}{m\asg a}{Q}$.

\textbf{Rama falsa} ($\neg(a>b)\equiv a\le b$): $\sust{Q}{b}{m}=b\ge a\land b\ge b\equiv b\ge a$. Como $a\le b\Rightarrow b\ge a$: $\hoare{\neg(a>b)}{m\asg b}{Q}$.

Por la regla condicional, la terna es válida (y total). Caso límite $a=b$: se ejecuta $m\asg b$ y $m=a=b$ cumple $Q$.
\end{sol}

\begin{prob}[\simil]
Demuestra $\hoare{a>b}{\kw{if}\ (a<b)\ \kw{then}\ m\asg b\ \kw{else}\ m\asg a}{m=a}$.
\end{prob}
\begin{sol}
$P: a>b$, $B: a<b$, $Q: m=a$.
\textbf{Rama verdadera:} $P\land B\equiv a>b\land a<b\equiv\F$: inalcanzable; $\hoare{\F}{m\asg b}{m=a}$ es válida por vacuidad (formalmente: $wp=(b=a)$ y $\F\Rightarrow b=a$).
\textbf{Rama falsa:} $P\land\neg B\equiv a>b\land a\ge b\equiv a>b$; $wp(m\asg a,\ m=a)=(a=a)\equiv\T$ y $a>b\Rightarrow\T$.
Por la regla condicional la terna es válida. $\blacksquare$
\end{sol}

\begin{prob}
¿Es válida $\hoare{\ }{\kw{if}\ a>b\ \kw{then}\ m\asg a\ \kw{else}\ m\asg b}{m>b}$? Si no, da un contraejemplo y propón dos formas de corregirla (cambiando solo la pre o solo la postcondición).
\end{prob}
\begin{sol}
\textbf{No es válida.} Rama falsa: $\neg(a>b)$ y $m\asg b$ exige $\sust{(m>b)}{b}{m}\equiv b>b\equiv\F$, pero $\neg(a>b)\not\Rightarrow\F$. Contraejemplo: $a=b=3$: $m=3$ y $3>3$ es falso.
Correcciones: (i) postcondición $m\ge b$ (válida en ambas ramas); (ii) precondición $a>b$: así la rama falsa se vuelve inalcanzable y la verdadera da $m=a>b$. (Con $a\neq b$ \emph{no} basta: si $a<b$, $m=b$ y $b>b$ es falso.)
\end{sol}

\begin{prob}
Demuestra $\hoare{\ }{\kw{if}\ x<0\ \kw{then}\ y\asg -x\ \kw{else}\ y\asg x}{y\ge0\land y^2=x^2}$ (valor absoluto).
\end{prob}
\begin{sol}
\textbf{Rama verdadera} ($x<0$): $\sust{Q}{-x}{y}=(-x\ge0)\land(x^2=x^2)\equiv x\le0$. Como $x<0\Rightarrow x\le0$: válida.
\textbf{Rama falsa} ($x\ge0$): $\sust{Q}{x}{y}=(x\ge0)\land(x^2=x^2)\equiv x\ge0$: coincide con la condición de la rama.
Regla condicional $\Rightarrow$ válida. Es decir, $y=|x|$.
\end{sol}

\begin{prob}
(\textbf{if} sin \textbf{else}) Demuestra $\hoare{\ }{\kw{if}\ x>y\ \kw{then}\ (t\asg x;\ x\asg y;\ y\asg t)}{x\le y}$.
\end{prob}
\begin{sol}
Regla: probar $\hoare{P\land B}{C_1}{Q}$ y $(P\land\neg B)\Rightarrow Q$.

\textbf{Rama verdadera}, hacia atrás desde $x\le y$:
por $y\asg t$: $x\le t$; \ por $x\asg y$: $y\le t$; \ por $t\asg x$: $y\le x$. Así $\hoare{y\le x}{t\asg x;\ x\asg y;\ y\asg t}{x\le y}$, y como $x>y\Rightarrow y\le x$, por consecuencia $\hoare{x>y}{C_1}{x\le y}$.

\textbf{Rama falsa (skip):} $\T\land\neg(x>y)\equiv x\le y\Rightarrow x\le y$. \checkmark

Por la regla del \textbf{if} sin \textbf{else}, la terna es válida.
\end{sol}

\begin{prob}
Demuestra $\hoare{x>0}{\kw{if}\ x>10\ \kw{then}\ x\asg x-10\ \kw{else}\ x\asg x+5}{x>0}$.
\end{prob}
\begin{sol}
\textbf{Rama verdadera:} $P\land B\equiv x>10$; $wp=x-10>0\iff x>10$. \checkmark
\textbf{Rama falsa:} $P\land\neg B\equiv 0<x\le10$; $wp=x+5>0\iff x>-5$, y $0<x\le10\Rightarrow x>-5$. \checkmark
Regla condicional $\Rightarrow$ válida. (Nótese que $x>0$ es un invariante de este bloque.)
\end{sol}

\begin{prob}[\reto]
Demuestra que el máximo de tres valores es correcto:
\[\kw{if}\ a\ge b\ \kw{then}\ (\kw{if}\ a\ge c\ \kw{then}\ m\asg a\ \kw{else}\ m\asg c)\ \kw{else}\ (\kw{if}\ b\ge c\ \kw{then}\ m\asg b\ \kw{else}\ m\asg c)\]
con $P=\T$ y $Q: m\ge a\land m\ge b\land m\ge c$.
\end{prob}
\begin{sol}
Se aplica la regla condicional tres veces; quedan 4 hojas, cada una con su precondición acumulada:
\begin{enumerate}
\item $a\ge b\land a\ge c$, $m\asg a$: $wp=a\ge a\land a\ge b\land a\ge c$, implicada. \checkmark
\item $a\ge b\land a<c$, $m\asg c$: $wp=c\ge a\land c\ge b\land c\ge c$; $c>a\ge b$ da $c\ge a$ y $c\ge b$. \checkmark
\item $a<b\land b\ge c$, $m\asg b$: $wp=b\ge a\land b\ge b\land b\ge c$. \checkmark
\item $a<b\land b<c$, $m\asg c$: $c>b>a$. \checkmark
\end{enumerate}
Cada \textbf{if} interno se demuestra con precondición $a\ge b$ (o $a<b$) y la misma $Q$; luego el \textbf{if} externo. Válida y total.
\end{sol}

% =====================================================================
\section{Ciclos: inducción, invariantes y terminación}
% =====================================================================

\begin{prob}[\examen]
Demostrar que el invariante en las siguientes instrucciones es $i+j=9$:
\begin{lstlisting}
int j = 9;
for( int i=0; i<10; i++ )
    j--;
\end{lstlisting}
¿Qué valores tienen $i$ y $j$ al terminar? ¿Por qué termina?
\end{prob}
\begin{sol}
Equivalente con \textbf{while}: \ \texttt{j:=9; i:=0; while i<10 do \{ j:=j-1; i:=i+1 \}}.

Sean $i_n$, $j_n$ los valores tras $n$ iteraciones. Traza: $(i,j)=(0,9),(1,8),(2,7),\dots$ \ Sea $P(n): i_n+j_n=9$.

\textbf{Caso base} ($n=0$, antes de entrar): $i_0=0$, $j_0=9\Rightarrow i_0+j_0=9$. \checkmark

\textbf{Hipótesis de inducción:} tras $n$ iteraciones, $i_n+j_n=9$.

\textbf{Paso inductivo:} si la guarda $i_n<10$ se cumple, el cuerpo hace $j_{n+1}=j_n-1$ e $i_{n+1}=i_n+1$. Entonces
\[i_{n+1}+j_{n+1}=(i_n+1)+(j_n-1)=i_n+j_n\overset{\text{H.I.}}{=}9.\]
Por inducción, $i+j=9$ vale después de cualquier número de iteraciones: es invariante. $\blacksquare$

\emph{(Forma Hoare equivalente:} $\hoare{i+j=9\land i<10}{j\asg j-1;\ i\asg i+1}{i+j=9}$. Hacia atrás: por $i\asg i+1$: $i+1+j=9$; por $j\asg j-1$: $i+1+j-1=9\iff i+j=9$. \checkmark)

\textbf{Salida:} también es invariante $0\le i\le10$ (empieza en 0, sube de uno en uno y solo sube si $i<10$). Al salir, $\neg(i<10)\land i\le10\Rightarrow i=10$, y por el invariante $j=9-10=\mathbf{-1}$.

\textbf{Terminación:} variante $V=10-i$: entero, $\ge0$ mientras el ciclo continúa y decrece en 1 cada iteración. Se ejecutan exactamente 10 iteraciones: corrección total.
\end{sol}

\begin{prob}[\simil]
Encuentra y demuestra el invariante que relaciona $k$ e $i$, y el valor final de $k$:
\begin{lstlisting}
int k = 20;
for( int i = 0; i < 5; i++ )
    k = k - 4;
\end{lstlisting}
\end{prob}
\begin{sol}
Traza: $(i,k)=(0,20),(1,16),(2,12),(3,8),(4,4),(5,0)$. Se observa $k_n=20-4n$, $i_n=n$; eliminando $n$: $\mathbf{I: k+4i=20}$.

Base: $k_0+4i_0=20+0=20$. \ Paso: $k_{n+1}+4i_{n+1}=(k_n-4)+4(i_n+1)=k_n+4i_n=20$. Por inducción, invariante.
Salida: $i=5\Rightarrow k=20-20=\mathbf 0$. Variante $5-i$.
\end{sol}

\begin{prob}
Demuestra que la función calcula el cuadrado de un entero $A\ge0$ (identifica invariante, salida y variante):
\begin{lstlisting}
Function Cuadrado( A )
  C (*$\leftarrow$*) 0;  D (*$\leftarrow$*) 0
  While ( D (*$\neq$*) A )
    C (*$\leftarrow$*) C + A
    D (*$\leftarrow$*) D + 1
  Return C
\end{lstlisting}
¿Qué ocurre si $A<0$?
\end{prob}
\begin{sol}
Traza con $A=4$: $(C,D)=(0,0),(4,1),(8,2),(12,3),(16,4)$: $C_n=nA$, $D_n=n\Rightarrow\mathbf{I: C=D\cdot A}$.

\textbf{Base:} $C_0=0=0\cdot A=D_0A$. \ \textbf{H.I.:} $C_n=D_nA$. \ \textbf{Paso:} $C_{n+1}=C_n+A=D_nA+A=(D_n+1)A=D_{n+1}A$. Por inducción, $I$ es invariante.

\textbf{Salida:} el ciclo termina con $D=A$; entonces $C=DA=A\cdot A=A^2$. (El invariante solo no basta: se combina con $\neg B$.)

\textbf{Terminación:} con $A\ge0$, $D$ parte de 0 y sube de 1 en 1; la variante $V=A-D$ es entera, $\ge0$ y decrece en 1: tras $A$ iteraciones $D=A$. Corrección total.

Si $A<0$: $D$ nunca alcanza $A$ (solo crece) y el ciclo \textbf{no termina}. La terna solo sería parcialmente correcta (``si terminara, $C=A^2$'').
\end{sol}

\begin{prob}
¿Qué realiza la subrutina? Demuéstralo ($N,M$ enteros, $M\ge0$).
\begin{lstlisting}
Subroutine ?( N, M; R )
  R (*$\leftarrow$*) 1
  While ( M > 0 )
    R (*$\leftarrow$*) R (*$\times$*) N
    M (*$\leftarrow$*) M (*$-$*) 1
  Return
End Subroutine
\end{lstlisting}
¿Por qué $R=N^M$ \emph{no} es un invariante?
\end{prob}
\begin{sol}
Calcula $R=N^{M_0}$, donde $M_0$ es el valor original de $M$.
Traza con $N=2$, $M_0=5$: $(R,M)=(1,5),(2,4),(4,3),(8,2),(16,1),(32,0)$; siempre $R\cdot2^M=32$.

\textbf{Invariante:} $I: R\cdot N^{M}=N^{M_0}$ (``calculado $\times$ pendiente $=$ total'').
\textbf{Base:} $R_0N^{M_0}=1\cdot N^{M_0}$. \checkmark\
\textbf{Paso:} si $M_n>0$: $R_{n+1}N^{M_{n+1}}=(R_nN)N^{M_n-1}=R_nN^{M_n}\overset{\text{H.I.}}=N^{M_0}$. \checkmark\
\textbf{Salida:} $M=0\Rightarrow R\cdot N^0=R=N^{M_0}$.
\textbf{Terminación:} variante $V=M$, entero $\ge0$ que decrece en 1.

$R=N^M$ no se conserva: al inicio $R=1$ y $N^M=N^{M_0}$ (falla si $M_0>0$, $N\ne1$). $R$ crece mientras $M$ baja: la relación conservada es el \emph{producto} $RN^M$. Error típico: expresar el resultado con el $M$ final (que vale 0) en lugar de $M_0$.
\end{sol}

\begin{prob}
Demuestra la corrección total de ($n\ge0$):
\begin{lstlisting}
s := 0; i := 0;
while i (*$\neq$*) n do
    i := i + 1;
    s := s + i
\end{lstlisting}
con postcondición $s=\frac{n(n+1)}{2}$.
\end{prob}
\begin{sol}
\textbf{Invariante:} $I: s=\frac{i(i+1)}{2}\land 0\le i\le n$.
\textbf{Base:} $i=0,s=0=\frac{0\cdot1}{2}$. \checkmark
\textbf{Paso:} si $i_n\ne n$ (y $i_n\le n$, luego $i_n<n$): $i_{n+1}=i_n+1\le n$ y
$s_{n+1}=s_n+i_{n+1}=\frac{i_n(i_n+1)}2+(i_n+1)=\frac{(i_n+1)(i_n+2)}{2}=\frac{i_{n+1}(i_{n+1}+1)}2$. \checkmark
\textbf{Salida:} $i=n\Rightarrow s=\frac{n(n+1)}2$. \quad \textbf{Variante:} $n-i$.
\end{sol}

\begin{prob}
División entera por restas sucesivas ($a\ge0$, $b>0$):
\begin{lstlisting}
q := 0; r := a;
while r (*$\geq$*) b do
    r := r - b;
    q := q + 1
\end{lstlisting}
Demuestra que al terminar $q=a\ \mathrm{div}\ b$ y $r=a\bmod b$. Traza con $a=17$, $b=5$.
\end{prob}
\begin{sol}
\textbf{Invariante:} $I: a=q\cdot b+r\land r\ge0$.
\textbf{Base:} $q=0$, $r=a$: $0\cdot b+a=a$ y $a\ge0$. \checkmark
\textbf{Paso:} con guarda $r\ge b$: $r'=r-b\ge0$, $q'=q+1$ y $q'b+r'=qb+b+r-b=qb+r=a$. \checkmark
\textbf{Salida:} $I\land\neg(r\ge b)\Rightarrow a=qb+r\land0\le r<b$, que es exactamente la definición de cociente y residuo.
\textbf{Terminación:} variante $r$: entera, $\ge0$ (por $I$) y decrece en $b\ge1$ cada iteración.
Traza: $(q,r)=(0,17),(1,12),(2,7),(3,2)$; sale con $r=2<5$: $17=3\cdot5+2$. \checkmark
\end{sol}

\begin{prob}
Para cada ciclo, propón el invariante y demuéstralo por inducción ($n\ge0$):
(a) \texttt{f:=1; i:=0; while i<n do \{ i:=i+1; f:=f*i \}} \qquad
(b) \texttt{s:=0; i:=0; while i<n do \{ s:=s+2*i+1; i:=i+1 \}}
\end{prob}
\begin{sol}
(a) $I: f=i!$. Base: $f=1=0!$. Paso: $f_{n+1}=f_n\cdot i_{n+1}=i_n!\,(i_n+1)=(i_n+1)!=i_{n+1}!$. Salida $i=n$: $f=n!$. Variante $n-i$.

(b) $I: s=i^2$ (suma de los primeros $i$ impares). Base: $0=0^2$. Paso: $s_{n+1}=s_n+2i_n+1=i_n^2+2i_n+1=(i_n+1)^2=i_{n+1}^2$. Salida: $s=n^2$. Variante $n-i$.
\end{sol}

\begin{prob}
Responde: (a) ¿Por qué una tabla de valores no es una demostración? (b) ¿Qué relación hay entre el caso base/paso inductivo y las obligaciones de Hoare para \textbf{while}? (c) ¿Qué se obtiene si se prueba el invariante y la salida, pero no la variante?
\end{prob}
\begin{sol}
(a) La tabla muestra algunos casos y sugiere el patrón; la inducción prueba que vale para \emph{cualquier} número de iteraciones.
(b) Caso base $\leftrightarrow$ inicialización ($P\Rightarrow I$); paso inductivo $\leftrightarrow$ mantenimiento $\hoare{I\land B}{C}{I}$; después se añade la salida $I\land\neg B\Rightarrow Q$.
(c) Solo corrección \emph{parcial}: ``si el ciclo termina, el resultado es correcto''. La variante (entera, no negativa, estrictamente decreciente) aporta la terminación y con ella la corrección total.
\end{sol}

% =====================================================================
\section{Función 91 de McCarthy, inducción modificada y Ackermann}
% =====================================================================

\begin{prob}
Con $f(x)=x-10$ si $x>100$ y $f(x)=f(f(x+11))$ si $x\le100$:
(a) desarrolla a mano $f(99)$ y $f(95)$ y cuenta cuántas llamadas se hacen en cada caso;
(b) desarrolla $f(89)$ mostrando solo los pasos intermedios importantes;
(c) calcula $f(105)$ y $f(101)$.
\end{prob}
\begin{sol}
(a) $f(99)=f(f(110))=f(100)=f(f(111))=f(101)=91$. Llamadas: $f(99),f(110),f(100),f(111),f(101)$: \textbf{5}.\\
$f(95)=f(f(106))=f(96)=f(f(107))=f(97)=\dots=f(100)=f(f(111))=f(101)=91$. Se llama a $f(95),\dots,f(101)$ (7) y a $f(106),\dots,f(111)$ (6): \textbf{13} llamadas.

(b) $f(89)=f(f(100))$; como $f(100)=91$ (3 llamadas), queda $f(91)=f(f(102))=f(92)=\dots=f(100)=f(101)=91$. (Total: 25 llamadas.)

(c) $f(105)=95$; \ $f(101)=91$. En general $f(x)=x-10$ para $x>100$ y $f(x)=91$ para $x\le101$.
\end{sol}

\begin{prob}
Enuncia el principio de inducción matemática modificado y úsalo para demostrar que $f(x)=91$ para todo entero $x\le100$. Explica por qué la inducción simple (``si vale en $n$ vale en $n+1$'') no basta.
\end{prob}
\begin{sol}
\textbf{Principio modificado.} Sea $k$ un entero fijo. Si (1) $P(k)$ es verdadera y (2) para todo $n<k$, [$P(m)$ verdadera para todo $m$ con $n<m\le k$] implica $P(n)$, entonces $P(n)$ vale para todo $n\le k$.

\textbf{Demostración} con $k=100$ y $P(x): f(x)=91$.

\emph{Paso básico:} $f(100)=f(f(111))=f(101)=91$ (pues $111>100$ y $101>100$).

\emph{Hipótesis:} sea $x<100$ y supongamos $f(m)=91$ para todo $m$ con $x<m\le100$.

\emph{Paso inductivo:} como $x\le100$, $f(x)=f(f(x+11))$.
\begin{itemize}
\item \textbf{Caso a)} $x+11>100$ (es decir, $90\le x\le99$): $f(x+11)=x+1$, luego $f(x)=f(x+1)$. Como $x<x+1\le100$, por H.I.\ $f(x+1)=91$.
\item \textbf{Caso b)} $x+11\le100$ ($x\le89$): como $x<x+11\le100$, por H.I.\ $f(x+11)=91$, luego $f(x)=f(91)$. Como $x\le89<91\le100$, por H.I.\ $f(91)=91$.
\end{itemize}
En ambos casos $f(x)=91$. Por el principio modificado, $f(x)=91$ para todo $x\le100$. $\blacksquare$

\textbf{Por qué no basta la inducción simple:} el caso b) necesita el valor de $f$ en \emph{dos} puntos mayores que $x$ ($x+11$ y $91$), no solo en $x+1$; la hipótesis debe cubrir todo el tramo $(x,100]$. Además la recursión ``se aleja'' (el argumento sube) antes de resolverse, así que no hay un contador que baje de uno en uno.
\end{sol}

\begin{prob}[\simil]
Sea $g(x)=x-2$ si $x>20$ y $g(x)=g(g(x+3))$ si $x\le20$. Calcula $g(20)$, $g(19)$, $g(15)$; conjetura el valor de $g(x)$ para $x\le20$ y demuéstralo con inducción modificada.
\end{prob}
\begin{sol}
$g(20)=g(g(23))=g(21)=19$; \ $g(19)=g(g(22))=g(20)=19$; \ $g(15)=g(g(18))=g(19)=19$ (usando $g(18)=g(g(21))=g(19)=19$).
\textbf{Conjetura:} $g(x)=19$ para todo $x\le20$.

\textbf{Paso básico} ($k=20$): $g(20)=19$ (calculado).
\textbf{H.I.:} para $x<20$, $g(m)=19$ para todo $m$ con $x<m\le20$.
\textbf{Paso:} $g(x)=g(g(x+3))$.
\begin{itemize}
\item Caso a) $x+3>20$ ($x\in\{18,19\}$): $g(x+3)=x+1$, así $g(x)=g(x+1)$ con $x<x+1\le20$: por H.I.\ vale 19.
\item Caso b) $x+3\le20$ ($x\le17$): $x<x+3\le20\Rightarrow g(x+3)=19$; luego $g(x)=g(19)$ y $x\le17<19\le20\Rightarrow g(19)=19$.
\end{itemize}
Por inducción modificada, $g(x)=19$ para todo $x\le20$. $\blacksquare$ \ (Patrón general del tipo McCarthy: techo $c$, resta $a$, suma $b>a$: valor $c+b-2a$; aquí $20+3-4=19$ y $100+11-20=91$.)
\end{sol}

\begin{prob}
Con la definición $A(0,y)=y+1$, $A(x+1,0)=A(x,1)$, $A(x+1,y+1)=A(x,A(x+1,y))$:
(a) calcula $A(1,2)$ y $A(2,1)$ aplicando \emph{solo} la definición;
(b) demuestra por inducción sobre $z$ que $A(1,z)=z+2$ y $A(2,z)=2z+3$.
\end{prob}
\begin{sol}
(a) $A(1,0)=A(0,1)=2$; \ $A(1,1)=A(0,A(1,0))=A(0,2)=3$; \ $A(1,2)=A(0,A(1,1))=A(0,3)=\mathbf 4$.\\
$A(2,0)=A(1,1)=3$; \ $A(2,1)=A(1,A(2,0))=A(1,3)=A(0,A(1,2))=A(0,4)=\mathbf 5$.

(b) \textbf{$A(1,z)=z+2$.} Base: $A(1,0)=A(0,1)=2$. H.I.: $A(1,n)=n+2$. Paso:
$A(1,n+1)=A(0,A(1,n))=A(0,n+2)=n+3=(n+1)+2$. \checkmark

\textbf{$A(2,z)=2z+3$.} Base: $A(2,0)=A(1,1)=3$. H.I.: $A(2,n)=2n+3$. Paso:
$A(2,n+1)=A(1,A(2,n))=A(1,2n+3)=(2n+3)+2=2(n+1)+3$, usando el resultado anterior como lema. \checkmark
\end{sol}

\begin{prob}
Demuestra que $A(3,z)=2^{z+3}-3$. Calcula $A(3,3)$ y $A(4,1)$. ¿Por qué el programa recursivo de Ackermann termina en \texttt{StackOverflowException} para valores pequeños como $(4,2)$?
\end{prob}
\begin{sol}
Base: $A(3,0)=A(2,1)=5=2^3-3$. Paso: $A(3,n+1)=A(2,A(3,n))=2(2^{n+3}-3)+3=2^{n+4}-6+3=2^{(n+1)+3}-3$. \checkmark

$A(3,3)=2^6-3=\mathbf{61}$. \ $A(4,0)=A(3,1)=13$; \ $A(4,1)=A(3,A(4,0))=A(3,13)=2^{16}-3=\mathbf{65533}$.
En general $A(4,z)$ es una torre de $z+3$ doses menos 3; $A(4,2)=2^{65536}-3$ tiene 19\,729 dígitos.

La recursión es doble y el argumento interno crece de forma no primitiva recursiva: la profundidad de la pila y el número de llamadas son astronómicos antes de llegar a un caso base; ninguna pila finita alcanza.
\end{sol}

\begin{prob}
Usando los lemas del problema anterior, demuestra $\hoare{x=2}{w\asg A(x,y)+1}{w=2y+4}$ y $\hoare{x=1}{w\asg A(x,y)}{w=y+2}$. ¿Por qué no basta el axioma de asignación?
\end{prob}
\begin{sol}
Primera: por asignación, $\sust{(w=2y+4)}{A(x,y)+1}{w}\equiv A(x,y)+1=2y+4\equiv A(x,y)=2y+3$, así $\hoare{A(x,y)=2y+3}{w\asg A(x,y)+1}{w=2y+4}$. Por el lema $A(2,z)=2z+3$: $x=2\Rightarrow A(x,y)=2y+3$. Por consecuencia:
\[\dfrac{x=2\Rightarrow A(x,y)=2y+3\qquad\hoare{A(x,y)=2y+3}{w\asg A(x,y)+1}{w=2y+4}}{\hoare{x=2}{w\asg A(x,y)+1}{w=2y+4}}.\]
Segunda: idéntica con $\hoare{A(x,y)=y+2}{w\asg A(x,y)}{w=y+2}$ y el lema $x=1\Rightarrow A(x,y)=y+2$.

El axioma solo produce la precondición $A(x,y)=\dots$; para llegar a $x=1$ o $x=2$ hace falta un \emph{teorema} sobre $A$ (demostrado por inducción aparte) y la regla de consecuencia. Así se reutilizan lemas sin volver a desenrollar la recursión.
\end{sol}

\begin{prob}
Compara la inducción usada para ciclos con la usada para $f$ (McCarthy) y Ackermann: dirección, hipótesis y qué se demuestra.
\end{prob}
\begin{sol}
\begin{tabular}{@{}p{3cm}p{5.5cm}p{6.3cm}@{}}
\toprule & Ciclos & McCarthy / Ackermann\\\midrule
Dirección & sube desde $n=0$ (iteraciones) & McCarthy: baja desde un techo $k$; Ackermann: inducción anidada sobre $x$ y, dentro, sobre $y$\\
Hipótesis & $P(n)$ para el $n$ anterior & $P(m)$ para \emph{todo} $m\in(x,k]$ (fuerte)\\
Qué prueba & un invariante se conserva & la función devuelve cierto valor/fórmula\\
Ejemplo & $C=DA$ en \texttt{Cuadrado} & $f(x)=91$; $A(2,z)=2z+3$\\\bottomrule
\end{tabular}
\end{sol}

% =====================================================================
\section{Programación lógica y Prolog}
% =====================================================================

\begin{prob}
(a) Escribe como cláusula (disyunción de literales) y como implicación la regla \texttt{abuelo(X,Z) :- padre(X,Y), padre(Y,Z).} (b) ¿Qué es una cláusula de Horn definida? ¿Y una cláusula objetivo? Representa la consulta \texttt{?- abuelo(juan, W).}
\end{prob}
\begin{sol}
(a) Implicación: $\forall X,Y,Z\;(padre(X,Y)\land padre(Y,Z)\to abuelo(X,Z))$. Cláusula: $abuelo(X,Z)\lor\neg padre(X,Y)\lor\neg padre(Y,Z)$.

(b) \textbf{Horn}: a lo más un literal positivo. \textbf{Definida}: exactamente uno (hechos y reglas; la cabeza es el positivo). \textbf{Objetivo}: ningún literal positivo; la consulta es $\leftarrow abuelo(juan,W)\equiv\neg abuelo(juan,W)$. Un hecho $padre(juan,pedro)$ es $padre(juan,pedro)\leftarrow\T$.
\end{sol}

\begin{prob}
Da el unificador más general o explica por qué no existe:
(a) $p(X,b)$ y $p(a,Y)$; \ (b) $f(X,g(Y))$ y $f(h(Z),g(a))$; \ (c) $p(X,X)$ y $p(a,b)$; \ (d) $gusta(maria,X)$ y $gusta(Y,cine)$; \ (e) $q(f(X),X)$ y $q(Y,a)$; \ (f) $p(X)$ y $p(f(X))$.
\end{prob}
\begin{sol}
(a) $\{X/a,\ Y/b\}$. \ (b) $\{X/h(Z),\ Y/a\}$. \ (c) No: $X/a$ obliga a $a=b$, constantes distintas. \ (d) $\{Y/maria,\ X/cine\}$. \ (e) $\{Y/f(a),\ X/a\}$ (tras $X/a$, $f(X)$ se vuelve $f(a)$). \ (f) No (prueba de ocurrencia): $X=f(X)$ exigiría un término infinito.
\end{sol}

\begin{prob}
Demuestra por \textbf{refutación} con resolución:
(a) $S=\{\neg hombre(X)\lor mortal(X),\ hombre(socrates)\}\models mortal(socrates)$.
(b) $S=\{P(X)\lor Q(X),\ \neg Q(X)\lor R(X),\ \neg R(X)\lor P(X)\}\models P(X)$.
\end{prob}
\begin{sol}
Se agrega la negación de la conclusión y se busca la cláusula vacía $\square$.

(a) $C_1:\neg hombre(X)\lor mortal(X)$, $C_2: hombre(socrates)$, $C_3:\neg mortal(socrates)$.
$C_3,C_1$ con $\theta=\{X/socrates\}$: $\neg hombre(socrates)$; con $C_2$: $\square$. Luego $S\models mortal(socrates)$.

(b) $C_4:\neg P(X)$.
$C_4,C_3\Rightarrow\neg R(X)$; \ $\neg R(X),C_2\Rightarrow\neg Q(X)$; \ $\neg Q(X),C_1\Rightarrow P(X)$; \ $P(X),C_4\Rightarrow\square$.
$S\cup\{\neg P(X)\}$ es insatisfactible $\Rightarrow S\models P(X)$.
\end{sol}

\begin{prob}
Dado el programa
\begin{lstlisting}[language=Prolog]
padre(juan, pedro).
padre(pedro, luis).
padre(pedro, ana).
abuelo(X, Z) :- padre(X, Y), padre(Y, Z).
\end{lstlisting}
(a) Escribe la derivación SLD para \texttt{?- abuelo(juan, W).} indicando en cada paso la cláusula usada y el unificador. (b) ¿Qué responde Prolog al teclear \texttt{;}?
\end{prob}
\begin{sol}
Renombramos la regla: $abuelo(X_1,Z_1)\leftarrow padre(X_1,Y_1),padre(Y_1,Z_1)$.
\begin{enumerate}
\item $G_0:\ \leftarrow abuelo(juan,W)$. Con la regla, $\theta_1=\{X_1/juan,\ Z_1/W\}$.
\item $G_1:\ \leftarrow padre(juan,Y_1),\ padre(Y_1,W)$. Primera meta con el hecho $padre(juan,pedro)$: $\theta_2=\{Y_1/pedro\}$.
\item $G_2:\ \leftarrow padre(pedro,W)$. Con $padre(pedro,luis)$: $\theta_3=\{W/luis\}$.
\item $G_3:\ \square$. Respuesta: \texttt{W = luis}.
\end{enumerate}
(b) Con \texttt{;} Prolog retrocede al último punto de elección ($G_2$) y usa $padre(pedro,ana)$: \texttt{W = ana}. Con otro \texttt{;} no quedan cláusulas para $padre(pedro,W)$ ni para $padre(juan,Y_1)$: \texttt{false}.
\end{sol}

\begin{prob}
Explica por qué el paso ``reemplazar una meta por el cuerpo de una regla'' es una instancia de la regla de resolución. ¿Qué significan las letras S, L, D?
\end{prob}
\begin{sol}
La regla $A\leftarrow B_1,\dots,B_m$ es la cláusula $A\lor\neg B_1\lor\dots\lor\neg B_m$ y la meta $\leftarrow A_1,\dots,A_k,\dots,A_n$ es $\neg A_1\lor\dots\lor\neg A_k\lor\dots\lor\neg A_n$. Si $\theta=\mathrm{UMG}(A,A_k)$, los literales $A$ y $\neg A_k$ son complementarios: se eliminan y el resolvente es
$(\neg A_1\lor\dots\lor\neg A_{k-1}\lor\neg B_1\lor\dots\lor\neg B_m\lor\neg A_{k+1}\lor\dots\lor\neg A_n)\theta$, que en notación de metas es $\leftarrow(A_1,..,A_{k-1},B_1,..,B_m,A_{k+1},..,A_n)\theta$.

\textbf{S}: regla de \emph{selección} de la meta (Prolog: la de más a la izquierda). \textbf{L}: \emph{lineal} (cada resolvente se obtiene del anterior y una cláusula del programa). \textbf{D}: cláusulas \emph{definidas}. Prolog además prueba las cláusulas en orden de programa, en profundidad, con retroceso.
\end{sol}

\begin{prob}
Dado el programa, indica si cada consulta tiene éxito, falla o produce error, y justifica:
\begin{lstlisting}[language=Prolog]
a.
b.
c :- fail.
s :- a, \+ c.
t :- c ; b.
u :- a, c.
v :- \+ a ; c.
w :- b, ( \+ c, a ).
x :- y.
\end{lstlisting}
Consultas: \texttt{?- s.} \texttt{?- t.} \texttt{?- u.} \texttt{?- v.} \texttt{?- w.} \texttt{?- x.}
\end{prob}
\begin{sol}
\begin{tabular}{@{}llp{11cm}@{}}
\texttt{s} & \checkmark & \texttt{a} es hecho; \texttt{c} no se puede demostrar (su cuerpo es \texttt{fail}), así que \texttt{\textbackslash+ c} tiene éxito.\\
\texttt{t} & \checkmark & disyunción: \texttt{c} falla, se prueba la alternativa \texttt{b}, que es hecho.\\
\texttt{u} & \texttimes & \texttt{a} tiene éxito pero \texttt{c} falla.\\
\texttt{v} & \texttimes & \texttt{\textbackslash+ a} falla (porque \texttt{a} sí se demuestra) y \texttt{c} también falla.\\
\texttt{w} & \checkmark & \texttt{b}; luego \texttt{\textbackslash+ c} tiene éxito y \texttt{a} es hecho.\\
\texttt{x} & error & \texttt{y} no está definido: SWI-Prolog lanza \emph{Unknown procedure: y/0} (\texttt{existence\_error}); no es lo mismo que ``falso''.
\end{tabular}

Recordatorio: \texttt{\textbackslash+ G} es \emph{negación como falla}: tiene éxito si el programa \emph{no logra demostrar} \texttt{G}.
\end{sol}

\begin{prob}
Con el programa \texttt{MariaGusta.pl}
\begin{lstlisting}[language=Prolog]
gusta( maria, helado ).
gusta( maria, leer ).
gusta( maria, cine ).
dar( juan, beso, maria ).
recibe( X, beso ) :- dar( juan, beso, X ).
\end{lstlisting}
indica la respuesta de: (a) \texttt{gusta(maria, cine).} (b) \texttt{gusta(juan, cine).} (c) \texttt{gusta(Quien, helado).} (d) \texttt{gusta(maria, Maria).} (con \texttt{;}) (e) \texttt{recibe(X, beso).} (f) \texttt{recibe ( X, beso ).}
\end{prob}
\begin{sol}
(a) \texttt{true} (átomos que coinciden). \ (b) \texttt{false}. \ (c) \texttt{Quien = maria}. \
(d) \texttt{Maria} es una \emph{variable} (mayúscula inicial), no el átomo \texttt{maria}: responde \texttt{Maria = helado ; Maria = leer ; Maria = cine}. \
(e) Se resuelve con la regla: la meta \texttt{dar(juan,beso,X)} unifica con el hecho: \texttt{X = maria}. \
(f) \textbf{Error de sintaxis} (\emph{operator expected}): no puede haber espacio entre el nombre del predicado y el paréntesis.
\end{sol}

\begin{prob}[\simil]
Considera
\begin{lstlisting}[language=Prolog]
maximo( X, Y, X ) :- X >= Y.
maximo( X, Y, Y ) :- X < Y.
mayor( X, Y, X ) :- X >= Y, !.
mayor( X, Y, Y ) :- X < Y.
max( X, Y, Y ) :- X < Y, !.
max( X, _, X ).
\end{lstlisting}
(a) ¿Qué hace el corte \texttt{!}? (b) ¿Qué es un corte ``verde'' y uno ``rojo''? Clasifica los de \texttt{mayor} y \texttt{max}. (c) ¿Qué responde \texttt{?- max(3, 5, 3).}? ¿Es correcto? Corrígelo.
\end{prob}
\begin{sol}
(a) El corte tiene éxito una vez y \emph{elimina los puntos de elección} creados desde que se eligió la cláusula actual: Prolog ya no probará otras cláusulas del mismo predicado ni otras alternativas de las metas anteriores al \texttt{!} dentro del cuerpo.

(b) \textbf{Verde:} solo mejora eficiencia; quitarlo no cambia las respuestas. \textbf{Rojo:} cambia el significado; quitarlo cambia las respuestas. En \texttt{mayor} es verde (las condiciones $X\ge Y$ y $X<Y$ son excluyentes; el corte evita probar una cláusula que igual fallaría). En \texttt{max} es rojo: la segunda cláusula no tiene condición y depende del corte (sin él, \texttt{max(3,5,M)} daría \texttt{M = 5 ; M = 3}).

(c) \texttt{?- max(3,5,3).} La primera cláusula exige unificar el tercer argumento con $Y=5$: $3\ne5$, falla \emph{antes} de llegar al corte. La segunda, \texttt{max(X,\_,X)}, unifica con $X=3$: responde \texttt{true}. \textbf{Incorrecto} (el máximo es 5).
Corrección: hacer la unificación de la salida \emph{después} del corte:
\begin{lstlisting}[language=Prolog]
max( X, Y, Z ) :- X < Y, !, Z = Y.
max( X, _, X ).
\end{lstlisting}
o usar la versión \texttt{maximo} sin cortes.
\end{sol}

\begin{prob}
Indica la respuesta de SWI-Prolog: (a) \texttt{X = 2+3.} (b) \texttt{X is 2+3.} (c) \texttt{5 is 2+3.} (d) \texttt{2+3 = 5.} (e) \texttt{2+3 =:= 5.} (f) \texttt{X is Y+1.}
\end{prob}
\begin{sol}
(a) \texttt{X = 2+3} (unifica con el \emph{término} sin evaluar). (b) \texttt{X = 5} (\texttt{is} evalúa el lado derecho). (c) \texttt{true}. (d) \texttt{false} (el término \texttt{+(2,3)} no es el número 5). (e) \texttt{true} (\texttt{=:=} compara valores aritméticos). (f) Error: \emph{arguments are not sufficiently instantiated} (\texttt{Y} no tiene valor).
\end{sol}

\begin{prob}
(a) Traza \texttt{?- factorial(3, F).} con
\begin{lstlisting}[language=Prolog]
factorial(0, 1) :- !.
factorial(N, F) :- N > 0, N1 is N - 1, factorial(N1, F1), F is N * F1.
\end{lstlisting}
(b) Escribe \texttt{suma\_hasta(N, S)} ($S=1+2+\dots+N$) y \texttt{mcd(A, B, G)} (algoritmo de Euclides) con el mismo estilo.
\end{prob}
\begin{sol}
(a) \texttt{factorial(3,F)}: $3>0$, $N1=2$, llama \texttt{factorial(2,F1)} $\to$ $N1=1$, \texttt{factorial(1,F1')} $\to$ $N1=0$, \texttt{factorial(0,F1'')}: primera cláusula, $F1''=1$ y el corte descarta la segunda. Al regresar: $F1'=1\cdot1=1$, $F1=2\cdot1=2$, $F=3\cdot2=\mathbf 6$.
\begin{lstlisting}[language=Prolog]
suma_hasta(0, 0) :- !.
suma_hasta(N, S) :- N > 0, N1 is N - 1, suma_hasta(N1, S1), S is S1 + N.

mcd(A, 0, A) :- !.
mcd(A, B, G) :- B > 0, R is A mod B, mcd(B, R, G).
\end{lstlisting}
\texttt{?- mcd(48, 18, G).} $\to$ \texttt{mcd(18,12)} $\to$ \texttt{mcd(12,6)} $\to$ \texttt{mcd(6,0)} $\to$ \texttt{G = 6}. Mismo patrón que en clase: caso base protegido por corte y caso recursivo que reduce el problema.
\end{sol}

\begin{prob}
Compara resolución general y resolución SLD (qué cláusulas usan, forma de la derivación, qué significa el éxito, cómo manejan alternativas). ¿Por qué el orden de las cláusulas importa en Prolog aunque lógicamente no importe?
\end{prob}
\begin{sol}
\begin{tabular}{@{}p{3.3cm}p{5.8cm}p{6cm}@{}}
\toprule & General & SLD / Prolog\\\midrule
Cláusulas & arbitrarias & definidas $+$ una meta\\
Elección & cualquier par de cláusulas & la meta seleccionada con una cláusula del programa\\
Derivación & árbol/grafo amplio & lineal por rama\\
Éxito & se deriva $\square$ & la lista de metas queda vacía\\
Alternativas & muchas combinaciones & retroceso (backtracking)\\\bottomrule
\end{tabular}

Lógicamente el conjunto de consecuencias no cambia al reordenar. Operacionalmente Prolog elige la meta de más a la izquierda y prueba cláusulas en orden, en profundidad: un orden desafortunado puede ser lento o \emph{no terminar} (p.\,ej.\ una regla recursiva por la izquierda como \texttt{anc(X,Y) :- anc(X,Z), padre(Z,Y).} antes del caso base). Además, que una búsqueda no termine no significa que la consulta sea falsa.
\end{sol}

% =====================================================================
\newpage
\section{Exámenes simulados}
% =====================================================================
Resuelve cada simulacro en \textbf{90 minutos}, sin notas, escribiendo las demostraciones completas como si fuera el examen real. Los problemas se numeran como \textbf{10.1--10.6} (simulacro A) y \textbf{10.7--10.12} (simulacro B).

\subsection*{Simulacro A}

\begin{prob}[ · A1]
Demostrar que la siguiente terna de Hoare es válida:
\[\hoare{a\neq b}{\kw{if}\ (a>b)\ \kw{then}\ (m\asg a;\ n\asg b)\ \kw{else}\ (m\asg b;\ n\asg a)}{m>n}.\]
\end{prob}
\begin{sol}
$P: a\ne b$, $B: a>b$, $Q: m>n$.

\textbf{Rama verdadera} $C_1: m\asg a;\ n\asg b$. Hacia atrás: por $n\asg b$: $m>b$; por $m\asg a$: $a>b$. Así $\hoare{a>b}{C_1}{m>n}$, y $P\land B\equiv a\ne b\land a>b\Rightarrow a>b$. \checkmark

\textbf{Rama falsa} $C_2: m\asg b;\ n\asg a$. Hacia atrás: por $n\asg a$: $m>a$; por $m\asg b$: $b>a$. Así $\hoare{b>a}{C_2}{m>n}$. Además $P\land\neg B\equiv a\ne b\land a\le b\equiv a<b\Rightarrow b>a$. \checkmark

Por la regla condicional (y composición en cada rama), la terna es válida; sin ciclos, corrección total. Nótese que la precondición $a\ne b$ es \emph{necesaria}: con $a=b$ se obtendría $m=n$.
\end{sol}

\begin{prob}[ · A2]
¿Qué es la corrección parcial y la corrección total de un programa? Da un ejemplo de programa parcialmente correcto que no sea totalmente correcto.
\end{prob}
\begin{sol}
Parcial: $\hoare{P}{C}{Q}$ es válida: si $C$ inicia en $P$ \emph{y termina}, termina en $Q$. Total: además $C$ termina para toda entrada en $P$. Total $=$ parcial $+$ terminación.

Ejemplo: $\hoare{\T}{\kw{while}\ x\ne0\ \kw{do}\ x\asg x-2}{x=0}$. Es parcialmente correcto: si el ciclo termina, es porque $x=0$ (invariante trivial $\T$, salida $\neg(x\ne0)\Rightarrow x=0$). No es totalmente correcto: con $x=3$ o $x=-2$ nunca termina. Con precondición $x\ge0\land x$ par, la variante $x$ da terminación. Otro ejemplo: \texttt{Cuadrado(A)} con $A<0$.
\end{sol}

\begin{prob}[ · A3]
Escribe en seudocódigo un algoritmo de divide y vencerás que calcule la suma de los elementos de $A[l..r]$. Plantea su recurrencia y traza la llamada sobre $A=[4,7,1,3]$.
\end{prob}
\begin{sol}
\begin{lstlisting}
Funcion Suma( A, l, r ): entero
Inicio
  Si l = r entonces devolver( A[l] )          // problema simple
  m = (l + r) div 2                           // descomponer en k = 2
  S1 = Suma( A, l, m )
  S2 = Suma( A, m+1, r )
  devolver( S1 + S2 )                         // combinar
Fin
\end{lstlisting}
$T(1)=\Theta(1)$, $T(n)=2T(n/2)+\Theta(1)\Rightarrow T(n)=\Theta(n)$.
Traza: \texttt{Suma(1,4)} $\to$ \texttt{Suma(1,2)} $=$ \texttt{Suma(1,1)}$+$\texttt{Suma(2,2)} $=4+7=11$; \ \texttt{Suma(3,4)} $=1+3=4$; \ resultado $11+4=15$.
\end{sol}

\begin{prob}[ · A4]
Demostrar que el invariante de las siguientes instrucciones es $s=2i$ y deducir el valor final de $s$ ($n\ge0$):
\begin{lstlisting}
int s = 0;
for( int i = 0; i < n; i++ )
    s = s + 2;
\end{lstlisting}
\end{prob}
\begin{sol}
Equivalente: \texttt{s:=0; i:=0; while i<n do \{ s:=s+2; i:=i+1 \}}.
\textbf{Base:} $s_0=0=2\cdot0=2i_0$. \ \textbf{H.I.:} $s_n=2i_n$. \ \textbf{Paso:} $s_{n+1}=s_n+2=2i_n+2=2(i_n+1)=2i_{n+1}$. Por inducción, $s=2i$ es invariante.
Además $0\le i\le n$ es invariante; al salir $i=n$ y por tanto $s=\mathbf{2n}$. Variante: $n-i$.
\end{sol}

\begin{prob}[ · A5]
Con el programa
\begin{lstlisting}[language=Prolog]
progenitor(ana, bea).
progenitor(ana, carlos).
progenitor(carlos, diana).
progenitor(bea, emilio).
abuela(X, Z) :- progenitor(X, Y), progenitor(Y, Z).
\end{lstlisting}
(a) Escribe la derivación SLD de \texttt{?- abuela(ana, N).} y todas las respuestas que da Prolog con \texttt{;}.
(b) ¿Qué responde \texttt{?- abuela(X, diana).}?
\end{prob}
\begin{sol}
(a) Regla renombrada: $abuela(X_1,Z_1)\leftarrow progenitor(X_1,Y_1),progenitor(Y_1,Z_1)$.
$G_0:\leftarrow abuela(ana,N)$; \ $\theta_1=\{X_1/ana,Z_1/N\}$.\\
$G_1:\leftarrow progenitor(ana,Y_1),progenitor(Y_1,N)$. Primer hecho que unifica: $progenitor(ana,bea)$, $\theta_2=\{Y_1/bea\}$.\\
$G_2:\leftarrow progenitor(bea,N)$; con $progenitor(bea,emilio)$: $\theta_3=\{N/emilio\}$; $G_3=\square$. \textbf{\texttt{N = emilio}}.\\
\texttt{;}: no hay otro $progenitor(bea,\_)$; se retrocede a $G_1$ y se usa $progenitor(ana,carlos)$: $Y_1/carlos$, $G_2':\leftarrow progenitor(carlos,N)$, $N/diana$: \textbf{\texttt{N = diana}}. Otro \texttt{;}: \texttt{false}.

(b) $G_1:\leftarrow progenitor(X,Y_1),progenitor(Y_1,diana)$. Se prueban los hechos en orden: $(ana,bea)\to progenitor(bea,diana)$ falla; $(ana,carlos)\to progenitor(carlos,diana)$ \checkmark: \texttt{X = ana}. Con \texttt{;} las demás ramas ($Y_1=diana$, $Y_1=emilio$) fallan: \texttt{false}.
\end{sol}

\begin{prob}[ · A6]
En el problema de las ocho reinas, explica para qué sirven $j-k$ y $j+k$, por qué basta el vector $\mathit{sol}[k]$ y qué parte del algoritmo realiza el ``backtracking''.
\end{prob}
\begin{sol}
Se coloca una reina por fila $k$ (así nunca comparten fila) y $\mathit{sol}[k]=j$ guarda su columna: esto describe completamente el tablero. Dos reinas se atacan si comparten columna ($j$), la diagonal donde $j-k$ es constante o la diagonal donde $j+k$ es constante; guardando esos tres identificadores ocupados, verificar una casilla no requiere recorrer el tablero (poda inmediata). El \emph{backtracking} es: probar columnas en orden, anotar la elegida (\texttt{push\_back}), bajar a la fila $k+1$ y, al regresar (por éxito o porque no hubo columna segura), deshacer (\texttt{pop\_back}) y probar la siguiente columna.
\end{sol}

\newpage
\subsection*{Simulacro B}

\begin{prob}[ · B1]
Encuentra la precondición más débil $P$ tal que $\hoare{P}{x\asg 2x+1;\ y\asg x-3}{y>4}$ sea válida. ¿Es válida $\hoare{x=5}{x\asg 2x+1;\ y\asg x-3}{y>4}$?
\end{prob}
\begin{sol}
Por $y\asg x-3$: $x-3>4\iff x>7$. \ Por $x\asg2x+1$: $2x+1>7\iff x>3$. \quad $P: x>3$.

Composición: $\hoare{x>3}{x\asg2x+1}{x>7}$ y $\hoare{x>7}{y\asg x-3}{y>4}$. Como $x=5\Rightarrow x>3$, por fortalecimiento de la precondición la segunda terna es \textbf{válida} (de hecho $x=11$, $y=8$).
\end{sol}

\begin{prob}[ · B2]
Describe el algoritmo de backtracking para encontrar \emph{todas} las soluciones de un problema. Aplícalo al problema de las 4 reinas: ¿cuáles son las soluciones?
\end{prob}
\begin{sol}
\begin{lstlisting}
Funcion Backtracking( Etapa i, solucionesCompletas ) devuelve: boolean
Inicio
  Éxito = falso;
  IniciarOpciones( i, GrupoOpciones o );
  Repetir
    SeleccionarNuevaOpcion( o, Opcion n );
    Si ( Aceptable( n ) ) entonces
      AnotarOpcion( i, n );
      Si SolucionCompleta( i ) entonces
        Éxito = verdadero;
        ApuntarSoluciónCompleta( i, solucionesCompletas );
      Sino
        Éxito = Backtracking( i+1 );
        Si Éxito = false entonces cancelamosAnotacion( i, n ); finsi;
      Finsi;
    Finsi;
  Hasta ( NoQuedanOpciones( o ) );
  Retorna Éxito;
Fin
\end{lstlisting}
A diferencia de ``una solución'', el ciclo no se detiene con el primer éxito: termina solo cuando no quedan opciones, y cada solución completa se agrega a la colección.

4 reinas (etapa $=$ fila, opción $=$ columna, aceptable $=$ columna y diagonales libres): soluciones $\mathit{sol}=[2,4,1,3]$ y $[3,1,4,2]$, es decir, casillas $(0,2),(1,4),(2,1),(3,3)$ y $(0,3),(1,1),(2,4),(3,2)$.
\end{sol}

\begin{prob}[ · B3]
Suponiendo demostrado que $A(1,z)=z+2$, demuestra que $A(2,z)=2z+3$ para todo $z\ge0$, calcula $A(2,5)$ y demuestra $\hoare{x=1}{w\asg A(x,y)}{w=y+2}$.
\end{prob}
\begin{sol}
Base: $A(2,0)=A(1,1)=3=2\cdot0+3$. H.I.: $A(2,n)=2n+3$. Paso: $A(2,n+1)=A(1,A(2,n))=A(1,2n+3)=2n+5=2(n+1)+3$. Por inducción, $A(2,z)=2z+3$. \ $A(2,5)=13$.

Terna: por asignación, $\hoare{A(x,y)=y+2}{w\asg A(x,y)}{w=y+2}$. Como $x=1\Rightarrow A(x,y)=A(1,y)=y+2$, por la regla de consecuencia $\hoare{x=1}{w\asg A(x,y)}{w=y+2}$. $\blacksquare$
\end{sol}

\begin{prob}[ · B4]
Para el ciclo de división entera por restas ($a\ge0$, $b>0$)
\begin{lstlisting}
q := 0; r := a;
while r (*$\geq$*) b do  { r := r - b; q := q + 1 }
\end{lstlisting}
encuentra el invariante, demuéstralo por inducción y prueba la corrección total respecto de $Q: a=qb+r\land0\le r<b$.
\end{prob}
\begin{sol}
Ver problema 7.6. Resumen: $I: a=qb+r\land r\ge0$. Base: $a=0\cdot b+a$, $a\ge0$. Paso: con $r_n\ge b$, $r_{n+1}=r_n-b\ge0$ y $q_{n+1}b+r_{n+1}=q_nb+b+r_n-b=a$. Salida: $I\land r<b\Rightarrow Q$. Terminación: variante $r\ge0$ que decrece en $b\ge1$.
\end{sol}

\begin{prob}[ · B5]
(a) ¿Qué responde Prolog a \texttt{?- max(3, 5, 3).} con \texttt{max(X,Y,Y) :- X < Y, !.} y \texttt{max(X,\_,X).}? Explica y corrige. (b) ¿Qué diferencia hay entre \texttt{X = 3*2}, \texttt{X is 3*2} y \texttt{3*2 =:= 6}?
\end{prob}
\begin{sol}
(a) Responde \texttt{true} (incorrecto): la primera cláusula falla al unificar el tercer argumento ($3$ con $Y=5$) antes de llegar al corte, y la segunda acepta cualquier caso. Es un \emph{corte rojo} mal usado. Corrección: \texttt{max(X,Y,Z) :- X < Y, !, Z = Y.} \ \texttt{max(X,\_,X).} (ver problema 9.8).

(b) \texttt{X = 3*2} liga \texttt{X} al \emph{término} \texttt{3*2} sin evaluar; \texttt{X is 3*2} evalúa y liga \texttt{X = 6}; \texttt{3*2 =:= 6} compara valores aritméticos: \texttt{true}.
\end{sol}

\begin{prob}[ · B6]
Demuestra con el principio de inducción modificado que la función 91 de McCarthy cumple $f(x)=91$ para todo entero $x\le100$.
\end{prob}
\begin{sol}
Ver la demostración completa del problema 8.2 (paso básico $f(100)=91$; hipótesis sobre todo $m\in(x,100]$; casos $x+11>100$ y $x+11\le100$). En el examen escribe los tres elementos explícitamente: \emph{principio usado, paso básico, hipótesis, paso inductivo por casos y conclusión}, verificando en cada uso de la hipótesis que el argumento cae en $(x,100]$.
\end{sol}

\vspace{1em}
\begin{consejo}[Lista de verificación antes del examen]
\begin{itemize}[leftmargin=*]
\item Sé escribir de memoria: esquema de divide y vencerás; backtracking para una, todas y la mejor solución.
\item Sé enunciar: axioma de asignación, consecuencia, composición, \textbf{if} con y sin \textbf{else}, obligaciones del \textbf{while}, inducción ordinaria y modificada.
\item Reconozco una rama inalcanzable ($P\land B\equiv\bot$) y sé justificarla.
\item Sé encontrar un invariante trazando 3--4 iteraciones y eliminando $n$; siempre incluyo salida y variante.
\item Sé derivar $A(1,z)$, $A(2,z)$, $A(3,z)$ y la prueba de $f(x)=91$.
\item Sé pasar reglas de Prolog a cláusulas, unificar, hacer una derivación SLD y predecir la consola (\texttt{;}, \texttt{!}, \texttt{\textbackslash+}, \texttt{is}, mayúsculas).
\item Sé definir cohesión, acoplamiento y métricas de eficiencia (tiempo, memoria).
\end{itemize}
\end{consejo}

